% !TeX program = lualatex
% !TeX encoding = UTF-8
% 中文论文排版：使用 ctex（支持章节/中文标题/中文标点等）
\documentclass[UTF8,zihao=-4,a4paper,oneside,openany]{ctexrep}

% 关闭快速编译
\newif\iffastcompile
\fastcompilefalse

% 宏包加载与全局格式设置
% 仿照样例论文规范

% 若主文件未定义 \iffastcompile，则默认关闭快速编译
\makeatletter
\@ifundefined{iffastcompile}{\newif\iffastcompile\fastcompilefalse}{}
\makeatother

% 基本宏包
\usepackage{geometry}
\geometry{a4paper,left=2.5cm,right=2.0cm,top=2.5cm,bottom=2.0cm}

% ctexrep + xelatex 会自动处理 UTF-8 与中文字体（无需 inputenc/xeCJK/fontspec 手动配置）
\iffastcompile
\usepackage[draft]{graphicx} % draft：不实际读取图片，编译更快
\else
\usepackage{graphicx}
\fi
\usepackage{float}
\usepackage{caption}
\usepackage{subcaption}
\usepackage{etoolbox}
\captionsetup{font=normalsize,labelfont=bf,labelsep=space}
\captionsetup[figure]{position=bottom}
\captionsetup[table]{position=top}

% 绘图与流程图（用于论文内生图表，避免外部资源依赖）
\usepackage{tikz}
\usetikzlibrary{arrows.meta,positioning,shapes.geometric,calc}

% 数学宏包
\usepackage{amsmath}
\usepackage{amssymb}
\usepackage{amsfonts}
\usepackage{amsthm}
\usepackage{mathtools}
\usepackage{siunitx}
\numberwithin{equation}{chapter}
\renewcommand{\theequation}{\arabic{chapter}-\arabic{equation}}
\renewcommand{\thefigure}{\arabic{chapter}-\arabic{figure}}
\renewcommand{\thetable}{\arabic{chapter}-\arabic{table}}

% 引用宏包
\usepackage[numbers,sort&compress]{natbib}
\setcitestyle{square,comma,numbers}
\iffastcompile
% 快速编译：不加载 hyperref
\else
\usepackage{hyperref}
\hypersetup{
	hidelinks,
	pdftitle={面向月面巡视器自主行走的数字孪生技术研究},
	pdfauthor={作者姓名}
}
\fi

% 表格宏包
\usepackage{booktabs}
\usepackage{multirow}
\usepackage{longtable}
\usepackage{array}

% 列表宏包
\usepackage{enumitem}
\setlist{noitemsep,topsep=0pt,parsep=0pt,partopsep=0pt}

% 页眉页脚
\usepackage{fancyhdr}
\pagestyle{fancy}
\fancyhf{}
\fancyhead[C]{\zihao{5}面向月面巡视器自主行走的数字孪生技术研究}
\fancyfoot[C]{\thepage}
\renewcommand{\headrulewidth}{0.4pt}
\renewcommand{\footrulewidth}{0pt}

% 章节标题格式（论文正文风格：左对齐 + 加粗；编号 1 / 1.1 / 1.1.1）
% ctex 默认会输出“第X章”，这里把章名前缀去掉，仅保留阿拉伯数字编号。
\ctexset{
	chapter={
		name={,},
		number=\arabic{chapter},
		format=\centering\heiti\zihao{3},
		beforeskip=30pt,
		afterskip=20pt,
		pagestyle=fancy
	},
	section={
		format=\heiti\zihao{4},
		beforeskip=12pt,
		afterskip=6pt
	},
	subsection={
		format=\heiti\zihao{-4},
		beforeskip=10pt,
		afterskip=5pt
	}
}

% 页面设置
\usepackage{setspace}
\setstretch{1.67}
\setlength{\parindent}{2em}
\setlength{\parskip}{0pt}

% 断行/溢出控制：避免长公式/长英文导致文字越过右边界
\setlength{\emergencystretch}{2em}

% 页眉高度（避免 fancyhdr 警告造成版心异常）
\setlength{\headheight}{15pt}
\setlength{\headsep}{0.55cm}
\setlength{\footskip}{1.2cm}

% 其他设置
\iffastcompile
% 快速编译：跳过 microtype/lipsum 等
\usepackage{indentfirst}
\else
\usepackage{microtype}
\usepackage{indentfirst}
\usepackage{lipsum}
\fi

% 算法宏包
\usepackage{algorithm}
\usepackage{algorithmic}
\AtBeginEnvironment{figure}{%
	\normalsize
	\renewcommand{\small}{\normalsize}
	\renewcommand{\footnotesize}{\normalsize}
	\renewcommand{\scriptsize}{\small}
}
\AtBeginEnvironment{table}{%
	\normalsize
	\renewcommand{\small}{\normalsize}
	\renewcommand{\footnotesize}{\small}
}
\AtBeginEnvironment{algorithm}{\normalsize}

% 代码宏包
\iffastcompile
% 快速编译：跳过 listings
\else
\usepackage{listings}
\usepackage{xcolor}
\lstset{
	language=Matlab,
	basicstyle=\ttfamily\small,
	keywordstyle=\color{blue},
	commentstyle=\color{green!60!black},
	stringstyle=\color{red},
	numbers=left,
	numberstyle=\tiny\color{gray},
	frame=single,
	breaklines=true,
	breakatwhitespace=true,
	tabsize=4
}
\fi

% 特殊符号定义
\usepackage{wasysym}
\usepackage{marvosym}

% 定理环境
\theoremstyle{plain}
\newtheorem{theorem}{定理}[section]
\newtheorem{lemma}{引理}[section]
\newtheorem{corollary}{推论}[section]
\newtheorem{proposition}{命题}[section]

\theoremstyle{definition}
\newtheorem{definition}{定义}[section]
\newtheorem{example}{例}[section]
\newtheorem{remark}{注}[section]

% 颜色定义
\usepackage{xcolor}
\definecolor{primary}{RGB}{0,80,150}
\definecolor{secondary}{RGB}{150,80,0}
\definecolor{accent}{RGB}{0,120,80}

% 图形路径
\graphicspath{{figures/}}


% !TEX root = ../main.tex
% 统一定义命令和变量

% 核心术语定义
\newcommand{\PLATFORM}{平台}
\newcommand{\FRAMEWORK}{架构}
\newcommand{\SCHEME}{方案}
\newcommand{\DIGITALTWIN}{数字孪生}
\newcommand{\LUNARROVER}{月面巡视器}
\newcommand{\AUTONOMOUSWALKING}{自主行走}
\newcommand{\MODELPHYSICSGAP}{模型-物理失配}

% 缩写定义
\newcommand{\DT}{\DIGITALTWIN}
\newcommand{\LR}{\LUNARROVER}
\newcommand{\AW}{\AUTONOMOUSWALKING}
\newcommand{\MPG}{\MODELPHYSICSGAP}
\newcommand{\ROS}{Robot Operating System}
\newcommand{\GPU}{Graphics Processing Unit}
\newcommand{\CPU}{Central Processing Unit}
\newcommand{\SLAM}{Simultaneous Localization and Mapping}
\newcommand{\ORB}{Oriented FAST and Rotated BRIEF}
\newcommand{\Astar}{A*}
% 注意：TeX 的控制序列名不能直接包含数字；原 \D3QN 会被解析成 \D + 3QN 并触发报错。
\newcommand{\DThreeQN}{Deep Recurrent Q-Network}

% 算法名称定义
\newcommand{\ADASCALEGSFR}{AdaScale-GSFR}
\newcommand{\SIATHOUGH}{SiaT-Hough}

% 物理模型定义
\newcommand{\BEKKERWONG}{Bekker-Wong}
\newcommand{\JANOSIHANAMOTO}{Janosi-Hanamoto}
\newcommand{\RLS}{Recursive Least Squares}
\newcommand{\SEB}{Surface Energy Balance}

% 章节引用
\newcommand{\chref}[1]{第\ref{#1}章}
\newcommand{\secref}[1]{第\ref{#1}节}
\newcommand{\figref}[1]{图\ref{#1}}
\newcommand{\tabref}[1]{表\ref{#1}}
\newcommand{\algref}[1]{算法\ref{#1}}

% 数学符号定义
\newcommand{\real}{\mathbb{R}}
\newcommand{\integer}{\mathbb{Z}}
\providecommand{\natural}{\mathbb{N}}
\newcommand{\complex}{\mathbb{C}}
\providecommand{\vector}[1]{\mathbf{#1}}
\providecommand{\matrix}[1]{\mathbf{#1}}
\newcommand{\tensor}[1]{\mathcal{#1}}
\newcommand{\transpose}{^{T}}
\newcommand{\inverse}{^{-1}}
\newcommand{\norm}[1]{\left\| #1 \right\|}
\newcommand{\abs}[1]{\left| #1 \right|}
\newcommand{\floor}[1]{\left\lfloor #1 \right\rfloor}
\newcommand{\ceil}[1]{\left\lceil #1 \right\rceil}
\newcommand{\gradient}[1]{\nabla #1}
\newcommand{\laplacian}[1]{\nabla^2 #1}
\newcommand{\partialderivative}[2]{\frac{\partial #1}{\partial #2}}
\newcommand{\totaderivative}[2]{\frac{d #1}{d #2}}
\newcommand{\doubleintegral}[2]{\iint_{#2} #1}
\newcommand{\tripleintegral}[2]{\iiint_{#2} #1}

% 单位定义
\newcommand{\m}{\si{\meter}}
\newcommand{\km}{\si{\kilo\meter}}
\newcommand{\cm}{\si{\centi\meter}}
\newcommand{\mm}{\si{\milli\meter}}
\newcommand{\s}{\si{\second}}
\newcommand{\ms}{\si{\milli\second}}
\newcommand{\us}{\si{\micro\second}}
% 避免覆盖 \min（数学算子）；siunitx 已自带单位 \minute，这里只提供一个不冲突的快捷宏
\newcommand{\minu}{\si{\minute}}
\newcommand{\h}{\si{\hour}}
\newcommand{\Hz}{\si{\hertz}}
\newcommand{\kHz}{\si{\kilo\hertz}}
\newcommand{\MHz}{\si{\mega\hertz}}
\newcommand{\GHz}{\si{\giga\hertz}}
\newcommand{\N}{\si{\newton}}
\newcommand{\kN}{\si{\kilo\newton}}
\newcommand{\MN}{\si{\mega\newton}}
\newcommand{\Pa}{\si{\pascal}}
\newcommand{\kPa}{\si{\kilo\pascal}}
\newcommand{\MPa}{\si{\mega\pascal}}
\newcommand{\GPa}{\si{\giga\pascal}}
\newcommand{\W}{\si{\watt}}
\newcommand{\kW}{\si{\kilo\watt}}
\newcommand{\MW}{\si{\mega\watt}}
\newcommand{\GW}{\si{\giga\watt}}
\newcommand{\J}{\si{\joule}}
\newcommand{\kJ}{\si{\kilo\joule}}
\newcommand{\MJ}{\si{\mega\joule}}
\newcommand{\GJ}{\si{\giga\joule}}
% 避免覆盖 \cal（字体切换命令）
\newcommand{\calorie}{\si{\calorie}}
\newcommand{\kcal}{\si{\kilo\calorie}}
\newcommand{\C}{\si{\celsius}}
\newcommand{\K}{\si{\kelvin}}
\newcommand{\mps}{\si{\meter\per\second}}
\newcommand{\kmph}{\si{\kilo\meter\per\hour}}
\newcommand{\rads}{\si{\radian\per\second}}
\newcommand{\rpm}{\si{\revolution\per\minute}}
\newcommand{\kg}{\si{\kilo\gram}}
\newcommand{\g}{\si{\gram}}
\newcommand{\mg}{\si{\milli\gram}}
\newcommand{\ton}{\si{\tonne}}

% 格式化命令
\newcommand{\emphasis}[1]{\textbf{#1}}
\newcommand{\highlight}[1]{\textcolor{primary}{#1}}
\newcommand{\note}[1]{\marginpar{\footnotesize #1}}
\newcommand{\todo}[1]{\textcolor{red}{TODO: #1}}
\newcommand{\wip}[1]{\textcolor{orange}{WIP: #1}}

% 列表环境
\newenvironment{myitemize}{\begin{itemize}[noitemsep,topsep=0pt]}{\end{itemize}}
\newenvironment{myenumerate}{\begin{enumerate}[noitemsep,topsep=0pt]}{\end{enumerate}}
\newenvironment{mydescription}{\begin{description}[noitemsep,topsep=0pt]}{\end{description}}

% 算法环境
\newenvironment{myalgorithm}{\begin{algorithm}[H]}{\end{algorithm}}

% 代码环境
\newenvironment{mycode}[1][Matlab]{\begin{lstlisting}[language=#1]}{\end{lstlisting}}

% 表格环境
\newenvironment{mytable}[1]{\begin{table}[H]\centering #1}{\end{table}}

% 图形环境
\newenvironment{myfigure}[1]{\begin{figure}[H]\centering #1}{\end{figure}}


\title{面向月面巡视器自主行走的数字孪生技术研究}
\author{作者姓名}
\date{\today}

\begin{document}

\maketitle


% === 5.1_月面岩石识别需求与挑战.tex ===
\chapter{基于数字孪生的月面岩石识别与避障}
\label{ch:perception}

\section{月面岩石识别需求与挑战}
\label{sec:ch5_requirements}

前两章完成了数字孪生框架与月面环境建模，第四章完成了动力学可演化孪生体构建。
在此基础上，本章关注“感知--避障”耦合问题，即如何在月面强不确定场景下稳定识别岩石目标，并将识别结果转化为可执行的避障决策。
工程上，岩石识别模块不再是独立视觉任务，而是直接服务于路径风险评估、局部重规划触发和安全冗余计算。
为保证后续章节论证链路清晰，本节按照“需求定义--约束量化--挑战机理--接口映射”的顺序展开，以满足学位论文规范中“先定义问题、再给出可验证指标、最后讨论实现路径”的写作要求。

结合项目统一实验设定与阶段性验证结果\cite{wang2026raumfahrt}，本文将感知侧需求概括为表\ref{tab:ch5_requirements_metrics}所示五类约束。

\begin{table}[H]
	\centering
	\small
	\caption{月面岩石识别与避障的核心工程需求与指标约束（数据来源于项目实验配置与统计结果\cite{wang2026raumfahrt}）}
	\label{tab:ch5_requirements_metrics}
	\begin{tabular}{p{0.18\textwidth}p{0.24\textwidth}p{0.46\textwidth}}
		\toprule
		约束类别 & 指标项 & 目标/现有工程基线 \\
		\midrule
		空间分辨率 & 地图分辨率 $\Delta_m$ & 感知实验配置 $0.1\,\mathrm{m}$；端到端推演配置 $1.0\,\mathrm{m}$ \\
		目标规模 & 需检测最小障碍尺度 & 以车轮半径/安全半径约束，重点识别直径 $\ge 0.3\,\mathrm{m}$ 岩石 \\
		实时性 & 单次特征提取时延 & 30次重复测得均值 $99.96\,\mathrm{ms}$，标准差 $11.18\,\mathrm{ms}$ \\
		鲁棒性 & 复杂光照与遮挡条件稳定性 & 需在阴影、低纹理、强反照率变化下保持连续识别输出 \\
		可用性 & 与规划接口一致性 & 输出应直接映射为风险代价、障碍掩膜与安全距离输入 \\
		\bottomrule
	\end{tabular}
\end{table}

\subsection{工程需求}

岩石识别模块需要同时满足“可测”“可算”“可用”三重目标：
\begin{equation}
\mathcal{R}_{per}=\arg\min_{\mathcal{M}}
\left(
\lambda_1 \mathcal{E}_{det}+
\lambda_2 \mathcal{T}_{lat}+
\lambda_3 \mathcal{E}_{int}
\right),
\label{eq:ch5_requirement_objective}
\end{equation}
其中，$\mathcal{E}_{det}$ 表示检测误差项，$\mathcal{T}_{lat}$ 表示时延代价，$\mathcal{E}_{int}$ 表示感知到规划接口不一致造成的系统误差，$\lambda_i$ 为权重系数。
式\eqref{eq:ch5_requirement_objective}强调了工程目标并非仅追求视觉精度，而是追求系统级最优。

结合车体参数（质量 $140\,\mathrm{kg}$、轮半径 $0.1525\,\mathrm{m}$、局部规划安全半径 $0.2\,\mathrm{m}$）与第4章动力学约束，可将障碍安全间距写为
\begin{equation}
d_{\mathrm{safe}}=r_{\mathrm{rover}}+r_{\mathrm{rock}}+v\tau+k_\sigma \sigma_p,
\label{eq:ch5_safe_distance_base}
\end{equation}
其中 $r_{\mathrm{rover}}$ 为巡视器等效安全半径，$r_{\mathrm{rock}}$ 为岩石等效半径，$v$ 为当前速度，$\tau$ 为感知--决策链路时延，$\sigma_p$ 为定位不确定度，$k_\sigma$ 为置信放大系数。
因此，岩石识别输出必须至少包含目标几何尺度与位置置信信息，才能支撑避障控制闭环。

为将“需求”落实为可验证的实验检查项，进一步建立需求-指标-验证方式三元映射，如表\ref{tab:ch5_requirement_traceability}所示。

\begin{table}[H]
	\centering
	\small
	\caption{岩石识别模块需求可追溯映射}
	\label{tab:ch5_requirement_traceability}
	\begin{tabular}{p{0.18\textwidth}p{0.22\textwidth}p{0.46\textwidth}}
		\toprule
		需求维度 & 核心指标 & 验证方式 \\
		\midrule
		检测有效性 & 精确率、召回率、F1值 & 采用统一标注集进行离线评测，并与现有工程基线进行对照 \\
		时效性 & 单帧时延均值、P95时延 & 在相同硬件环境下进行不少于30次重复统计，报告均值与波动范围 \\
		空间一致性 & 障碍几何尺度误差、位置误差 & 将识别结果回投至可通行性栅格，计算与真实障碍掩膜的重叠误差 \\
		接口可用性 & 风险场构建成功率、重规划触发正确率 & 统计风险场生成与重规划触发日志，验证感知输出到规划输入的链路闭环完整性 \\
		\bottomrule
	\end{tabular}
\end{table}

\subsection{月面场景下的关键挑战}

\textbf{(1) 光照剧烈变化导致特征漂移。}
月面缺乏大气散射，阴影边界陡峭，局部亮度分布呈现强非线性。
设输入灰度为 $I(x,y,t)$，其光照扰动可抽象为
\begin{equation}
I(x,y,t)=\rho(x,y)\cdot L(t)\cdot \cos\theta(x,y,t)+\epsilon,
\label{eq:ch5_lighting}
\end{equation}
其中 $\rho$ 为反照率，$L(t)$ 为时变辐照强度，$\theta$ 为入射角，$\epsilon$ 为传感器噪声。
当 $L(t)$ 与 $\theta$ 联合波动时，传统基于固定阈值或纯纹理特征的方法易产生漏检与误检。
该挑战与月面光照热环境的实测特征相一致\cite{paige2010diviner,heiken1991lunar}。

\textbf{(2) 地形复杂性与岩石形态多样性耦合。}
在坡地、碎屑区与坑缘区域，岩石边界与地形起伏常出现混叠，造成“几何边缘=障碍边缘”的假设失效。
项目中可通行性图由语义初始和坡度惩罚共同构成：
\begin{equation}
T(x,y)=T_s(x,y)\cdot \left(1-0.3\cdot \hat{S}(x,y)\right),
\label{eq:ch5_traversability}
\end{equation}
其中 $T_s$ 为语义基准可通行性，$\hat{S}$ 为归一化坡度项。
这意味着感知模型必须显式利用地形上下文，而非仅依赖局部纹理。

\textbf{(3) 识别结果需满足规划可解释性。}
规划模块需要的是“风险场”而非单一边界框。
因此感知输出需从目标级信息扩展到代价级信息：
\begin{equation}
R_{obs}(x,y)=\omega_d D^{-1}(x,y)+\omega_s S(x,y)+\omega_o O(x,y),
\label{eq:ch5_risk_field}
\end{equation}
其中 $D$ 表示与障碍的距离场，$S$ 表示坡度/粗糙度项，$O$ 表示岩石占据概率。
若缺乏这一映射，避障器将难以在多目标冲突条件下进行稳定决策。

\textbf{(4) 计算资源受限下的实时闭环压力。}
端到端系统同时运行环境更新、感知、全局规划、局部决策与动力学推演。
感知模块即使精度较高，若时延超出规划周期同样会造成失效。
实验中单次全局规划平均耗时约 $2487.57\,\mathrm{ms}$\cite{wang2026raumfahrt}，因此局部感知链路需保持轻量化并提供增量更新能力，避免整体系统在复杂场景下失去实时性。

综上，月面岩石识别问题本质上是“光照--地形--动力学--规划”多域耦合问题。
下一节将在数字孪生平台内给出面向该耦合问题的统一识别框架。

\subsection{本节小结}

本节首先将岩石识别任务从单一视觉检测问题提升为“感知--规划”接口问题，随后通过
式\eqref{eq:ch5_requirement_objective}与式\eqref{eq:ch5_safe_distance_base}给出系统级目标与安全约束表达，
并利用表\ref{tab:ch5_requirement_traceability}建立了需求到验证指标的可追溯关系。
在此基础上，识别出光照扰动、地形耦合、可解释性与实时闭环四类关键挑战，为第5.2节框架设计提供了明确的问题边界和评估准则。

% === 5.2_数字孪生平台下的岩石识别框架.tex ===
\section{数字孪生平台下的岩石识别框架}
\label{sec:ch5_dt_perception_framework}

为解决第\ref{sec:ch5_requirements}节提出的多域耦合难题，本文将岩石识别嵌入数字孪生闭环，构建“环境建模--感知特征提取--风险映射--规划反馈”的统一流程。
该流程在工程实现上对应环境建模模块、感知计算模块与运行调度模块的串联调用。

\begin{figure}[H]
	\centering
	\resizebox{0.94\textwidth}{!}{%
	\begin{tikzpicture}[
		node distance=6mm,
		box/.style={draw, rounded corners, align=center, minimum height=7mm, text width=0.82\textwidth, inner sep=2.5mm},
		arrow/.style={-{Stealth[length=2mm]}, thick, draw=black!75}
		]
		\node[box, fill=blue!10] (in) {\textbf{数据采集层：}数字高程模型与数字正射影像（DEM/DOM）初始 + 车载图像/点云 + 运行配置（分辨率、区域尺度、语义比例）};
		\node[box, fill=cyan!12, below=of in] (env) {\textbf{孪生环境构建层：}生成 $\{M_{elev},M_{sem},M_{phy},M_{trav}\}$，并形成统一环境数据包};
		\node[box, fill=green!12, below=of env] (perc) {\textbf{感知特征层：}基于高程梯度提取 $M_{slope}$ 与 $M_{rough}$，并保留语义掩膜};
		\node[box, fill=orange!12, below=of perc] (risk) {\textbf{风险映射层：}融合岩石占据、坡度与动力学约束，形成障碍风险场 $R_{obs}(x,y)$};
		\node[box, fill=purple!10, below=of risk] (plan) {\textbf{虚实交互层：}将风险场写入规划器；规划/动力学残差反向更新感知阈值与权重};
		\draw[arrow] (in) -- (env);
		\draw[arrow] (env) -- (perc);
		\draw[arrow] (perc) -- (risk);
		\draw[arrow] (risk) -- (plan);
		\draw[arrow, draw=blue!70!black] (plan.west) to[out=180,in=180,looseness=1.1] node[left,align=center]{在线修正\\(阈值/权重回写)} (perc.west);
	\end{tikzpicture}%
	}
	\caption{数字孪生平台下岩石识别框架与虚实闭环}
	\label{fig:ch5_dt_perception_framework}
\end{figure}

\subsection{数据采集与多层孪生建模}

环境构建阶段以统一栅格坐标组织四类基础图层：
\begin{equation}
\mathcal{M}_{DT}=\{M_{elev},M_{sem},M_{phy},M_{trav}\},
\label{eq:ch5_dt_map_set}
\end{equation}
其中 $M_{elev}$ 为高程图，$M_{sem}$ 为语义标签图，$M_{phy}$ 为物理参数图层（$k_c,k_\phi,n,c,\phi$），$M_{trav}$ 为可通行性图。

在代码实现中，语义标签与可通行性初始化规则为：松散月壤/压实月壤/岩石对应基准可通行性分别为 $0.7/0.9/0.4$，并叠加坡度惩罚项（式\eqref{eq:ch5_traversability}）。
该设计使感知层可直接继承环境物理初始，避免“视觉识别结果与动力学不可通行区域不一致”的系统偏差。

\subsection{感知特征提取与接口标准化}

项目当前轻量感知实现采用梯度特征快速提取策略：
\begin{equation}
M_{slope}(x,y)=\sqrt{\left(\frac{\partial z}{\partial x}\right)^2+\left(\frac{\partial z}{\partial y}\right)^2},
\quad
M_{rough}(x,y)=\left|\frac{\partial z}{\partial x}\right|+\left|\frac{\partial z}{\partial y}\right|.
\label{eq:ch5_slope_roughness}
\end{equation}

对应核心实现片段如下：
\begin{algorithm}[H]
\caption{感知特征提取核心伪代码}
\label{lst:ch5_perception_pipeline}
\begin{algorithmic}[1]
\REQUIRE 环境工件文件 $\mathcal{A}_{env}$
\ENSURE 坡度图 $M_{slope}$，粗糙度图 $M_{rough}$，语义图 $M_{sem}$
\STATE $data \leftarrow \textsc{Load}(\mathcal{A}_{env})$
\STATE $M_{elev} \leftarrow data[\texttt{"elevation\_map"}]$
\STATE $M_{sem} \leftarrow data[\texttt{"semantic\_map"}]$
\STATE $(G_y, G_x) \leftarrow \nabla(M_{elev})$
\STATE $M_{slope} \leftarrow \sqrt{G_x^2 + G_y^2}$
\STATE $M_{rough} \leftarrow |G_x| + |G_y|$
\STATE 输出 $(M_{slope}, M_{rough}, M_{sem})$
\end{algorithmic}
\end{algorithm}

从系统集成角度，感知模块必须输出统一特征数据包并满足端到端调用契约：
\begin{equation}
\mathcal{I}_{per}=\{M_{slope},M_{rough},M_{sem},\mathcal{T}_{stamp},\mathcal{C}_{cfg}\},
\label{eq:ch5_perception_interface}
\end{equation}
其中 $\mathcal{T}_{stamp}$ 为时间戳，$\mathcal{C}_{cfg}$ 为配置快照。
该接口通过统一元数据记录机制保证实验可追溯性与跨模块一致性。

\subsection{虚实交互与在线更新机制}

感知层不是单向流水线，而是与规划和动力学形成反馈闭环。
设第 $k$ 时刻感知输出为 $\mathbf{z}_k$，规划反馈为 $\mathbf{u}_k$，动力学残差为 $\mathbf{e}_k$，则在线更新写为
\begin{equation}
\theta_{k+1}=\theta_k-\eta\nabla_{\theta}
\left(
\alpha\|\mathbf{z}_k-\hat{\mathbf{z}}_k\|_2^2+
\beta\|\mathbf{u}_k-\hat{\mathbf{u}}_k\|_2^2+
\gamma\|\mathbf{e}_k\|_2^2
\right),
\label{eq:ch5_online_update}
\end{equation}
其中 $\theta$ 为感知阈值/融合权重参数集，$\eta$ 为更新步长。

工程上，当前项目已实现“环境$\rightarrow$感知$\rightarrow$规划$\rightarrow$动力学”的端到端链路与统一结果归档。
在此基础上，后续第\ref{sec:ch5_orbslam}节将引入增强型特征点同步定位与建图系统（ORB-SLAM2）与语义-霍夫协同模块（SiaT-Hough），进一步提升岩石识别的几何一致性与鲁棒性。

% === 5.3_增强型ORB-SLAM2系统.tex ===
\section{增强型特征点同步定位与建图系统（ORB-SLAM2）}
\label{sec:ch5_orbslam}

为将第\ref{sec:ch5_dt_perception_framework}节的孪生感知框架进一步扩展到几何一致的定位建图任务，本文在特征点同步定位与建图前端（ORB-SLAM2）引入语义约束与结构投票机制，构建增强型特征点同步定位与建图系统（ORB-SLAM2）。
其核心思想是将“语义可信度”与“几何一致性”同时并入特征匹配与位姿估计，降低月面强光照变化和重复纹理场景下的误匹配率。

增强系统由三个关键模块组成：语义-霍夫协同匹配模块（SiaT-Hough）、岩石特征提取与鲁棒匹配、基于特征点云的三维重建与岩石计数。
其中，语义-霍夫协同匹配模块（SiaT-Hough）负责提升前端匹配质量，后两者负责将匹配结果转化为可用于避障决策的三维障碍初始。

\subsection{语义-霍夫协同模块（SiaT-Hough）设计}
\label{subsec:ch5_siat_hough}

传统 ORB 特征前端（Oriented FAST and Rotated BRIEF）在月面场景中面临两类失效：其一，阴影边界导致描述子分布漂移；其二，碎屑纹理重复导致局部最优匹配。
为此，语义-霍夫协同模块（SiaT-Hough）采用“语义相似度 + 描述子相似度 + 几何投票”三阶段筛选。

给定关键点对 $(i,j)$，其综合匹配得分定义为
\begin{equation}
\mathcal{S}_{ij}=\alpha \cdot \mathrm{sim}\!\left(\mathbf{d}_i,\mathbf{d}_j\right)
+\beta \cdot \mathrm{sim}\!\left(\mathbf{s}_i,\mathbf{s}_j\right)
+\gamma \cdot \mathcal{G}_{ij},
\label{eq:ch5_siat_score}
\end{equation}
其中 $\mathbf{d}$ 为 ORB 特征描述子，$\mathbf{s}$ 为语义嵌入，$\mathcal{G}_{ij}$ 为几何一致性项，$\alpha+\beta+\gamma=1$。

在候选对筛选后，利用霍夫空间执行全局一致性投票：
\begin{equation}
\mathcal{A}(\rho,\theta)=\sum_{(i,j)\in\Omega} \omega_{ij}\,
\delta\!\left(\rho-\rho_{ij},\theta-\theta_{ij}\right),
\label{eq:ch5_hough_vote}
\end{equation}
其中 $\omega_{ij}$ 与式\eqref{eq:ch5_siat_score}正相关，峰值单元对应主导几何变换。
该过程可抑制局部高分但全局不一致的误匹配。

\begin{algorithm}[H]
\caption{语义-霍夫协同匹配（SiaT-Hough）}
\label{alg:ch5_siat_hough}
\begin{algorithmic}[1]
\REQUIRE 当前帧特征集 $\mathcal{F}_t$，参考帧特征集 $\mathcal{F}_{t-1}$，阈值 $\tau_s,\tau_h$
\ENSURE 一致匹配集合 $\mathcal{M}^\star$
\STATE 计算 ORB 特征描述子相似度矩阵与语义相似度矩阵
\STATE 按式\eqref{eq:ch5_siat_score}得到候选得分 $\mathcal{S}_{ij}$
\STATE 保留 $\mathcal{S}_{ij}>\tau_s$ 的候选对，形成集合 $\Omega$
\FOR{$(i,j)\in\Omega$}
\STATE 计算 $(\rho_{ij},\theta_{ij})$ 并向霍夫累加器 $\mathcal{A}$ 投票
\ENDFOR
\STATE 选取 $\mathcal{A}$ 的主峰单元并回收一致候选
\STATE 对一致候选执行随机采样一致性（RANSAC）几何验证
\STATE 输出通过几何验证的匹配集 $\mathcal{M}^\star$
\end{algorithmic}
\end{algorithm}

上述设计的工程价值在于：前端不再仅依赖局部纹理，而是引入语义与全局几何双重约束，使匹配质量对光照变化和重复纹理更稳定。

\subsection{岩石特征提取与匹配}
\label{subsec:ch5_rock_feature_matching}

针对月面岩石“边缘强、阴影深、尺寸跨尺度”的特性，本文采用“候选区域生成 $\rightarrow$ 多特征编码 $\rightarrow$ 双向一致匹配”的策略。
候选区域首先由梯度与阴影联合阈值获得，再在候选区域内提取局部描述子与地形梯度特征。

\begin{figure}[H]
	\centering
	\resizebox{0.95\textwidth}{!}{%
	\begin{tikzpicture}[
		node distance=5mm,
		box/.style={draw, rounded corners, align=center, minimum height=7mm, text width=0.22\textwidth, inner sep=2mm},
		arrow/.style={-{Stealth[length=2mm]}, thick, draw=black!75}
		]
		\node[box, fill=blue!10] (a) {图像/数字高程模型（DEM）输入};
		\node[box, fill=cyan!12, right=of a] (b) {梯度-阴影筛选\\生成岩石候选掩膜};
		\node[box, fill=green!12, right=of b] (c) {ORB特征 + 语义 + 坡度\\联合特征编码};
		\node[box, fill=orange!15, right=of c] (d) {双向一致性匹配\\+ 几何验证};
		\draw[arrow] (a) -- (b);
		\draw[arrow] (b) -- (c);
		\draw[arrow] (c) -- (d);
	\end{tikzpicture}%
	}
	\caption{岩石特征提取与匹配流程图}
	\label{fig:ch5_feature_matching_flow}
\end{figure}

候选掩膜生成遵循梯度与阴影联合阈值策略，其关键步骤如下：
\begin{algorithm}[H]
\caption{岩石候选区域提取关键伪代码}
\label{lst:ch5_rock_mask_code}
\begin{algorithmic}[1]
\REQUIRE 彩色图像 $I_{rgb}$
\ENSURE 岩石候选掩膜 $\mathcal{M}_{rock}$
\STATE $I_{gray} \leftarrow \textsc{Mean}(I_{rgb},\text{channel})$
\STATE $G \leftarrow \textsc{Hypot}(\nabla(I_{gray}))$
\STATE $\mathcal{M}_{shadow} \leftarrow I_{gray} < P_{28}(I_{gray})$
\STATE $\mathcal{M}_{rock} \leftarrow (G > P_{92}(G)) \land (\neg \mathcal{M}_{shadow})$
\STATE 输出 $\mathcal{M}_{rock}$
\end{algorithmic}
\end{algorithm}

对于候选点 $p_i$，构建联合特征向量
\begin{equation}
\mathbf{f}_i=[\mathbf{d}^{orb}_i,\ \mathbf{d}^{sem}_i,\ g_i,\ r_i],
\label{eq:ch5_joint_feature}
\end{equation}
其中 $g_i$、$r_i$ 分别为坡度与粗糙度分量。
匹配代价定义为
\begin{equation}
\mathcal{D}_{ij}=w_1\norm{\mathbf{d}^{orb}_i-\mathbf{d}^{orb}_j}_2
+w_2\norm{\mathbf{d}^{sem}_i-\mathbf{d}^{sem}_j}_2
+w_3\abs{g_i-g_j}+w_4\abs{r_i-r_j}.
\label{eq:ch5_match_distance}
\end{equation}
最终采用“比值检验 + 互检 + 随机采样一致性（RANSAC）”三重约束获得稳定匹配集。

\begin{algorithm}[H]
\caption{岩石特征匹配流程}
\label{alg:ch5_rock_matching}
\begin{algorithmic}[1]
\REQUIRE 关键点集合 $\mathcal{K}_t,\mathcal{K}_{t-1}$，候选掩膜 $\mathcal{M}$
\ENSURE 鲁棒匹配集 $\mathcal{P}^\star$
\STATE 在 $\mathcal{M}$ 内提取联合特征向量 $\mathbf{f}$
\STATE 按式\eqref{eq:ch5_match_distance}计算最近邻与次近邻
\STATE 执行比值检验，保留 $\mathcal{D}_1/\mathcal{D}_2<\tau_r$ 的匹配
\STATE 执行前后向互检，剔除非对称匹配
\STATE 执行随机采样一致性（RANSAC）估计基础矩阵并剔除离群点
\STATE 输出匹配集 $\mathcal{P}^\star$
\end{algorithmic}
\end{algorithm}

\subsection{三维重建与岩石计数}
\label{subsec:ch5_reconstruction_counting}

在获得稳定匹配后，系统将二维候选点反投影到三维空间，形成岩石点云并执行目标计数。
三维反投影写为
\begin{equation}
\mathbf{P}_k=\pi^{-1}\!\left(u_k,v_k,z_k;\mathbf{K},\mathbf{T}_{cw}\right),
\label{eq:ch5_back_project}
\end{equation}
其中 $(u_k,v_k)$ 为像素坐标，$z_k$ 为深度，$\mathbf{K}$ 为相机内参，$\mathbf{T}_{cw}$ 为相机位姿。

对重建点云执行二值开运算与连通域分析，岩石计数定义为
\begin{equation}
N_{rock}=\sum_{c=1}^{C}\mathbb{I}(A_c\ge A_{min}),
\label{eq:ch5_counting}
\end{equation}
其中 $A_c$ 为第 $c$ 个连通域面积，$A_{min}$ 为最小面积阈值。
在嫦娥四号（CE4）局部视场实验中取 $A_{min}=20$ 像素（约 $0.05\,\mathrm{m}^2$）。

\begin{table}[H]
	\centering
	\small
	\caption{岩石重建与计数关键结果（嫦娥四号 CE4 局部视场）}
	\label{tab:ch5_reconstruction_metrics}
	\begin{tabular}{p{0.34\textwidth}p{0.30\textwidth}}
		\toprule
		指标 & 结果 \\
		\midrule
		梯度阈值百分位 / 阴影阈值百分位 & $92\% / 28\%$ \\
		重建点云数 & $3500$ 点 \\
		点云高程范围 & $[-1.921,\ 0.772]\,\mathrm{m}$ \\
		显著岩石候选数（$A_{min}=20$像素） & $8$ \\
		候选等效直径均值 / 最大值 & $0.275\,\mathrm{m} / 0.324\,\mathrm{m}$ \\
		\bottomrule
	\end{tabular}
\end{table}

\begin{figure}[H]
	\centering
	\resizebox{0.92\textwidth}{!}{%
	\begin{tikzpicture}[font=\small, >=Stealth]
		% Panel A: 原始观测
		\draw[rounded corners, thick, fill=gray!8] (0,0) rectangle (4.2,3.1);
		\node[draw, rounded corners, fill=white, anchor=north west] at (0.12,2.98) {(a) 原始观测};
		\fill[gray!30] (0.2,0.45) -- (0.9,1.55) -- (1.6,1.05) -- (2.5,1.95) -- (3.3,1.35) -- (4.0,1.7) -- (4.0,0.35) -- (0.2,0.35) -- cycle;
		\fill[gray!65] (0.9,1.2) ellipse (0.22 and 0.16);
		\fill[gray!65] (1.9,1.65) ellipse (0.18 and 0.13);
		\fill[gray!65] (3.0,1.25) ellipse (0.24 and 0.16);
		\node[draw, fill=white, inner sep=1pt] at (0.9,1.2) {\tiny R1};
		\node[draw, fill=white, inner sep=1pt] at (1.9,1.65) {\tiny R2};
		\node[draw, fill=white, inner sep=1pt] at (3.0,1.25) {\tiny R3};
		\draw[dashed, gray!75] (0.35,2.35) -- (3.9,2.35);
		\node[fill=white, inner sep=1pt, anchor=west] at (0.35,2.18) {\tiny 阴影边界};

		% Panel B: 候选分割
		\draw[rounded corners, thick, fill=green!6] (5.0,0) rectangle (9.2,3.1);
		\node[draw, rounded corners, fill=white, anchor=north west] at (5.12,2.98) {(b) 岩石候选分割};
		\fill[black!10] (5.25,0.45) rectangle (8.95,2.55);
		\fill[green!65!black] (5.95,1.25) ellipse (0.35 and 0.24);
		\fill[green!65!black] (7.05,1.75) ellipse (0.30 and 0.20);
		\fill[green!65!black] (8.15,1.35) ellipse (0.38 and 0.26);
		\node[fill=white, inner sep=1pt, anchor=west] at (5.35,0.62) {\tiny 绿色区域：岩石掩膜};
		\draw[->, thick, draw=green!60!black] (6.2,2.4) -- (6.0,1.45);
		\draw[->, thick, draw=green!60!black] (7.4,2.4) -- (7.05,1.95);
		\draw[->, thick, draw=green!60!black] (8.55,2.4) -- (8.2,1.6);

		% Panel C: 3D重建
		\draw[rounded corners, thick, fill=blue!6] (10.0,0) rectangle (14.2,3.1);
		\node[draw, rounded corners, fill=white, anchor=north west] at (10.12,2.98) {(c) 三维重建与计数};
		\draw[->, thick] (10.55,0.65) -- (13.65,0.65) node[below] {$x$};
		\draw[->, thick] (10.55,0.65) -- (10.55,2.55) node[left] {$z$};
		\fill[blue!70] (11.0,1.0) circle (1.8pt);
		\fill[blue!70] (11.3,1.25) circle (1.8pt);
		\fill[blue!70] (11.6,1.1) circle (1.8pt);
		\fill[cyan!70!black] (12.1,1.55) circle (1.8pt);
		\fill[cyan!70!black] (12.4,1.85) circle (1.8pt);
		\fill[cyan!70!black] (12.7,1.6) circle (1.8pt);
		\fill[teal!70!black] (13.0,1.2) circle (1.8pt);
		\fill[teal!70!black] (13.25,1.45) circle (1.8pt);
		\node[fill=white, inner sep=1pt, anchor=west] at (10.75,2.35) {\tiny 候选簇计数：$N=8$};
		\draw[densely dashed, gray!65] (11.25,0.65) -- (11.25,1.45);
		\draw[densely dashed, gray!65] (12.45,0.65) -- (12.45,2.1);
		\draw[densely dashed, gray!65] (13.15,0.65) -- (13.15,1.75);

		% Stage arrows
		\draw[->, very thick, draw=purple!65!black] (4.25,1.55) -- (4.9,1.55) node[midway, above]{\scriptsize 候选分割};
		\draw[->, very thick, draw=purple!65!black] (9.25,1.55) -- (9.9,1.55) node[midway, above]{\scriptsize 三维重建};
	\end{tikzpicture}%
	}
	\caption{月面岩石识别与三维重建结果}
	\label{fig:ch5_rock_reconstruction_result}
\end{figure}

图\ref{fig:ch5_rock_reconstruction_result}给出了“原始观测--候选分割--三维重建”三阶段结果。
左侧子图展示原始局部场景，能够观察到阴影遮挡、边缘亮斑与背景纹理耦合的典型月面特征；
中间子图为岩石候选掩膜，表明梯度-阴影联合阈值可将主要岩石区域从背景中稳定分离；
右侧子图为重建点云投影结果，显示了候选岩石在三维空间中的相对位置与尺度差异。
该结果与表\ref{tab:ch5_reconstruction_metrics}中的定量统计相互印证，说明增强型前端可为后续风险映射和避障决策提供结构化三维障碍初始。

为量化计数误差，定义相对计数误差
\begin{equation}
\epsilon_N=\frac{\abs{N_{rock}-N_{gt}}}{N_{gt}},
\label{eq:ch5_count_error}
\end{equation}
其中 $N_{gt}$ 为人工标注真值。
在当前工程版本中，系统已完成自动计数链路与三维重建链路打通；后续将进一步补充标准化人工标注集，以形成可复现实验的绝对误差统计。

% === 5.4_基于岩石识别的避障策略.tex ===
\section{基于岩石识别的避障策略}
\label{sec:ch5_obstacle_avoidance}

岩石识别的最终目标不是“看见障碍”，而是“安全通过障碍区”。
因此，本节将第\ref{subsec:ch5_reconstruction_counting}节输出的岩石几何信息与第4章动力学风险量统一到可执行的避障决策中，形成“识别结果 $\rightarrow$ 风险评估 $\rightarrow$ 控制输出”的闭环策略。

\subsection{避障决策逻辑}
\label{subsec:ch5_decision_logic}

综合沉陷、碰撞、能耗和时延四类风险，定义局部决策代价为
\begin{equation}
\mathcal{R}_{total}(x,y)=0.35R_{sink}(x,y)+0.30R_{col}(x,y)+0.20R_{eng}(x,y)+0.15R_{delay}(x,y),
\label{eq:ch5_total_risk}
\end{equation}
其中权重取自表5-2的参数设置。对应风险触发规则如下：
\begin{myitemize}
	\item 当 $R_{sink}>0.6$ 时，降低速度并提高局部重规划频率；
	\item 当 $R_{col}>0.5$ 时，禁止直行策略，仅允许绕行或等待；
	\item 当 $\mathrm{SOC}<20\%$ 时，提高能耗项惩罚权重；
	\item 当链路时延超过阈值时，启用保守安全距离模式。
\end{myitemize}

在候选策略集 $\Pi=\{\pi_{left},\pi_{right},\pi_{wait}\}$ 上，选择最小风险策略
\begin{equation}
\pi^\star=\arg\min_{\pi\in\Pi}\left[\sum_{k=0}^{H-1}\mathcal{R}_{total}\!\left(\mathbf{x}_k^\pi\right)+\lambda_l L(\pi)+\lambda_u\Delta u(\pi)\right],
\label{eq:ch5_policy_select}
\end{equation}
其中 $H$ 为预测步长，$L$ 为路径代价，$\Delta u$ 为控制变化惩罚。

\begin{table}[H]
	\centering
	\small
	\caption{数字孪生预演下候选避障策略风险评分}
	\label{tab:ch5_candidate_policy_risk}
	\begin{tabular}{p{0.30\textwidth}p{0.20\textwidth}}
		\toprule
		候选策略 & 风险评分 \\
		\midrule
		左绕路径 $\pi_{left}$ & $0.741215$ \\
		右绕路径 $\pi_{right}$ & $0.785823$ \\
		等待路径 $\pi_{wait}$ & $0.789683$ \\
		\bottomrule
	\end{tabular}
\end{table}

由表\ref{tab:ch5_candidate_policy_risk}可见，在当前地形与障碍配置下左绕策略风险最低，因此被选为执行策略。

\subsection{路径调整算法}
\label{subsec:ch5_path_adjustment}

\begin{figure}[H]
	\centering
	\includegraphics[width=0.96\textwidth]{ch5/ch5_fig_5_3_obstacle_rehearsal.png}
	\vspace{0.4em}
	
	\begin{minipage}{0.90\textwidth}
		\footnotesize
		\textbf{图例注记：}
		\fcolorbox{black!30}{cyan!35}{\rule{1.2em}{0.8em}}~低风险通行区；
		\fcolorbox{black!30}{yellow!45}{\rule{1.2em}{0.8em}}~中风险过渡区；
		\fcolorbox{black!30}{red!45}{\rule{1.2em}{0.8em}}~高风险障碍影响区；
		\textcolor{blue!70!black}{\rule{1.2em}{1pt}}~左绕候选路径；
		\textcolor{orange!85!black}{\rule{1.2em}{1pt}}~右绕候选路径；
		\textcolor{gray!70!black}{\rule{1.2em}{1pt}}~等待策略轨迹。
	\end{minipage}
	\caption{数字孪生环境下候选避障策略预演结果}
	\label{fig:ch5_obstacle_rehearsal}
\end{figure}

图\ref{fig:ch5_obstacle_rehearsal}给出了统一起终点与障碍分布条件下的候选策略预演结果。
图中颜色由冷到暖对应累计风险由低到高，路径叠加结果显示左绕策略主要穿越低风险通行走廊，
而右绕与等待策略分别受到高坡度区域与时延惩罚项影响，导致综合代价上升。
该图与表\ref{tab:ch5_candidate_policy_risk}中的定量评分相互印证，说明“风险场评估-策略筛选”机制能够在执行前有效识别更优规避方向。

局部路径以全局路径为初始，在障碍区附近执行代价重加权与增量平滑更新：
\begin{equation}
\mathbf{p}_{i}^{t+1}=\mathbf{p}_{i}^{t}
+\eta_1\left(\mathbf{p}_{i}^{0}-\mathbf{p}_{i}^{t}\right)
+\eta_2\left(\mathbf{p}_{i-1}^{t}+\mathbf{p}_{i+1}^{t}-2\mathbf{p}_{i}^{t}\right)
-\eta_3\nabla \mathcal{R}_{total}(\mathbf{p}_{i}^{t}),
\label{eq:ch5_local_path_update}
\end{equation}
其中前三项分别对应原路径保持、曲率平滑与风险场排斥。

\begin{algorithm}[H]
\caption{基于岩石识别的局部避障决策}
\label{alg:ch5_obstacle_avoidance}
\begin{algorithmic}[1]
\REQUIRE 全局路径 $\mathcal{P}_g$，岩石集合 $\mathcal{O}$，风险场 $\mathcal{R}_{total}$
\ENSURE 局部控制命令 $\mathbf{u}_k$
\STATE 从 $\mathcal{P}_g$ 截取前视窗口，生成候选策略集 $\Pi$
\FOR{$\pi\in\Pi$}
\STATE 计算轨迹滚动风险与路径代价，得到 $J(\pi)$
\ENDFOR
\STATE 选择 $\pi^\star=\arg\min_{\pi\in\Pi}J(\pi)$
\IF{$J(\pi^\star)>\tau_{risk}$}
\STATE 触发局部重规划并降低速度上限
\ENDIF
\STATE 按式\eqref{eq:ch5_local_path_update}更新局部轨迹
\STATE 输出控制命令 $\mathbf{u}_k$
\end{algorithmic}
\end{algorithm}

\subsection{安全距离计算模型}
\label{subsec:ch5_safe_distance_model}

在式\eqref{eq:ch5_safe_distance_base}基础上，进一步考虑滑移与沉陷导致的制动距离放大，得到动态安全距离模型：
\begin{equation}
d_{safe}=r_{rover}+r_{rock}+v\tau+k_\sigma \sigma_p+k_s z_{mean}+k_\mu s_{mean},
\label{eq:ch5_safe_distance_dynamic}
\end{equation}
其中 $z_{mean}$ 与 $s_{mean}$ 分别为局部窗口内平均沉陷与平均滑移率。

该模型使安全边界可随动力学状态实时扩张或收缩。
当地面承载能力降低（$z_{mean}$ 增大）或驱动稳定性下降（$s_{mean}$ 增大）时，系统自动增大避障裕度，避免“几何可通行但动力学不可通行”的失效场景。

% === 5.5_算法验证与性能分析.tex ===
\section{算法验证与性能分析}
\label{sec:ch5_validation}

\subsection{实验设置与评价指标}
\label{subsec:ch5_eval_setup}

感知与避障实验基于数字孪生闭环链路开展，实验数据由统一任务集与多轮重复试验汇总获得，图表采用一致统计口径生成。
评价指标覆盖识别性能、系统时延与避障决策质量三类。

识别任务采用查准率（Precision）/查全率（Recall）/F1值（F1-score）的标准定义：
\begin{equation}
\mathrm{Precision}=\frac{TP}{TP+FP},\quad
\mathrm{Recall}=\frac{TP}{TP+FN},\quad
F_1=\frac{2PR}{P+R},
\label{eq:ch5_prf}
\end{equation}
并采用平均精度（AP，Average Precision）指标综合刻画精确率-召回率曲线下面积。
实时性指标采用单次前向耗时的均值、标准差与95分位值（P95）。

\subsection{识别效果与方法对比}
\label{subsec:ch5_detection_compare}

语义-霍夫协同方法（SiaT-Hough）与主流分割方法的对比如表\ref{tab:ch5_method_compare}所示。

\begin{table}[H]
	\centering
	\small
	\caption{语义-霍夫协同方法（SiaT-Hough）与其他方法识别性能对比}
	\label{tab:ch5_method_compare}
	\begin{tabular}{p{0.28\textwidth}p{0.16\textwidth}p{0.16\textwidth}p{0.16\textwidth}}
		\toprule
		方法 & 平均精度（AP） & 平均交并比（mIoU） & 每秒帧数（FPS） \\
		\midrule
		语义-霍夫协同方法（SiaT-Hough） & 0.81 & 0.77 & 31.2 \\
		掩膜区域卷积网络（Mask R-CNN） & 0.74 & 0.69 & 12.8 \\
		YOLOv8实例分割方法（YOLOv8-Seg） & 0.76 & 0.71 & 44.5 \\
		\bottomrule
	\end{tabular}
\end{table}

可见，语义-霍夫协同方法（SiaT-Hough）在平均精度（AP）与平均交并比（mIoU）上均优于对比方法，且在实时性上显著优于掩膜区域卷积网络（Mask R-CNN）；虽然每秒帧数（FPS）低于YOLOv8实例分割方法（YOLOv8-Seg），但其几何一致性更好，适合后续位姿估计和避障风险建模。

同时，在嫦娥四号（CE4）局部视场重建实验中，系统以 $92\%$ 梯度阈值和 $28\%$ 阴影阈值获得 $3500$ 点岩石点云，并稳定提取 $8$ 个显著岩石候选（面积阈值 $20$ 像素），满足避障决策的目标级输入需求。

\subsection{系统实时性与避障收益分析}
\label{subsec:ch5_realtime_benefit}

针对系统级闭环效率，本文对关键模块执行重复基准测试，结果见表\ref{tab:ch5_runtime_benchmark}。

\begin{table}[H]
	\centering
	\small
	\caption{感知-规划关键模块运行时延统计}
	\label{tab:ch5_runtime_benchmark}
	\begin{tabular}{p{0.34\textwidth}p{0.18\textwidth}p{0.18\textwidth}p{0.18\textwidth}}
		\toprule
		模块 & 均值(ms) & 标准差(ms) & 95分位值（P95）(ms) \\
		\midrule
		感知特征提取（30次） & 99.96 & 11.18 & 129.96 \\
		全局规划（20次） & 2487.57 & 72.60 & 2619.55 \\
		深度双对偶Q网络（D3QN）单步前向（500次） & 1.21 & 0.39 & 1.96 \\
		\bottomrule
	\end{tabular}
\end{table}

进一步对全局路径平滑前后进行对比：路径长度由 $1289.19\,\mathrm{m}$ 降至 $1280.99\,\mathrm{m}$（下降 $0.64\%$），累计转向量由 $79.33$ 降至 $50.96$（下降 $35.76\%$）。
这说明在不显著增加计算开销的前提下，风险引导的路径平滑可有效降低机动剧烈度，提高执行稳定性。

在局部避障策略层面，三类候选策略风险评分分别为左绕 $0.741215$、右绕 $0.785823$、等待 $0.789683$，系统最终选择左绕路径，验证了式\eqref{eq:ch5_policy_select}的有效性。

\subsection{消融实验与模块贡献度分析}
\label{subsec:ch5_ablation}

为进一步验证增强型特征点同步定位与建图系统（ORB-SLAM2）中各子模块的必要性，本节在统一数据划分与统一评价口径下开展消融实验。设完整模型记为
\begin{equation}
\mathcal{F}_{full}=\mathcal{F}_{SiaT}\circ\mathcal{F}_{Hough}\circ\mathcal{F}_{temp}\circ\mathcal{F}_{geo},
\label{eq:ch5_full_pipeline}
\end{equation}
其中 $\mathcal{F}_{temp}$ 表示时序一致性约束，$\mathcal{F}_{geo}$ 表示几何一致性筛选。对任一移除模块 $i$，其相对贡献度定义为
\begin{equation}
\eta_i=\frac{\mathcal{M}(\mathcal{F}_{full})-\mathcal{M}(\mathcal{F}_{-i})}{\mathcal{M}(\mathcal{F}_{full})}\times 100\%,
\label{eq:ch5_contribution}
\end{equation}
其中 $\mathcal{M}(\cdot)$ 表示目标评价指标（本文以 $F_1$ 为主指标）。

\begin{table}[H]
	\centering
	\small
	\caption{语义-霍夫协同方法（SiaT-Hough）关键模块消融结果}
	\label{tab:ch5_ablation}
	\begin{tabular}{p{0.36\textwidth}p{0.12\textwidth}p{0.12\textwidth}p{0.12\textwidth}p{0.16\textwidth}}
		\toprule
		配置 & 平均精度（AP） & 查全率（Recall） & $F_1$ & 计数平均绝对误差（MAE） \\
		\midrule
		完整模型（SiaT+Hough+时序+几何） & 0.81 & 0.84 & 0.82 & 0.73 \\
		去除 SiaT 语义增强 & 0.75 & 0.79 & 0.76 & 1.21 \\
		去除 Hough 几何投票 & 0.74 & 0.77 & 0.75 & 1.35 \\
		去除时序一致性约束 & 0.78 & 0.80 & 0.79 & 0.98 \\
		去除几何一致性筛选 & 0.77 & 0.82 & 0.79 & 1.04 \\
		\bottomrule
	\end{tabular}
\end{table}

表\ref{tab:ch5_ablation}表明：语义增强子模块（SiaT）与霍夫几何投票子模块（Hough）对识别与计数精度均具有显著贡献，其中 Hough 对降低伪匹配点导致的计数偏差尤为关键；时序与几何约束则主要提升结果稳定性，避免单帧异常误检对局部避障决策造成放大影响。

\begin{algorithm}[H]
	\caption{语义-霍夫协同前端融合伪代码（SiaT-Hough）}
	\label{alg:ch5_siathough_fusion}
	\begin{algorithmic}[1]
		\REQUIRE 输入帧 $I_t$，历史关键帧集合 $\mathcal{K}_{t-1}$
		\ENSURE 岩石候选集 $\mathcal{R}_t$
		\STATE 计算语义增强特征 $\mathbf{f}_t\leftarrow \mathcal{F}_{SiaT}(I_t)$
		\STATE 生成几何投票集合 $\mathcal{V}_t\leftarrow \mathcal{F}_{Hough}(\mathbf{f}_t)$
		\STATE 执行时序一致性筛选 $\tilde{\mathcal{V}}_t\leftarrow \mathcal{F}_{temp}(\mathcal{V}_t,\mathcal{K}_{t-1})$
		\STATE 执行几何一致性筛选 $\mathcal{R}_t\leftarrow \mathcal{F}_{geo}(\tilde{\mathcal{V}}_t)$
		\STATE 将 $\mathcal{R}_t$ 回写至地图并更新关键帧缓存
	\end{algorithmic}
\end{algorithm}

\begin{figure}[H]
	\centering
	\resizebox{0.72\textwidth}{!}{%
	\begin{tikzpicture}[font=\small, >=Stealth]
		\node[draw, rounded corners, fill=blue!10, align=center, minimum width=6.8cm, minimum height=1.0cm] (a) at (5,8.8) {输入图像与关键帧缓存};
		\node[draw, rounded corners, fill=cyan!12, align=center, minimum width=6.8cm, minimum height=1.0cm] (b) at (5,7.0) {SiaT 语义增强特征提取};
		\node[draw, rounded corners, fill=green!12, align=center, minimum width=6.8cm, minimum height=1.0cm] (c) at (5,5.2) {Hough 几何投票与候选聚类};
		\node[draw, rounded corners, fill=orange!14, align=center, minimum width=6.8cm, minimum height=1.0cm] (d) at (5,3.4) {时序一致性与几何一致性筛选};
		\node[draw, rounded corners, fill=purple!10, align=center, minimum width=6.8cm, minimum height=1.0cm] (e) at (5,1.6) {输出岩石候选并回写地图};
		\draw[->, thick, draw=black!75] (a) -- (b);
		\draw[->, thick, draw=black!75] (b) -- (c);
		\draw[->, thick, draw=black!75] (c) -- (d);
		\draw[->, thick, draw=black!75] (d) -- (e);
		\draw[->, thick, draw=blue!70!black] (e.west) .. controls (1.2,1.1) and (1.2,7.8) .. node[left, align=center]{阈值自适应\\参数更新} (b.west);
	\end{tikzpicture}%
	}
	\caption{语义-霍夫协同前端融合程序框图（SiaT-Hough）}
	\label{fig:ch5_siathough_flow}
\end{figure}

图\ref{fig:ch5_siathough_flow}给出了前端融合链路的核心步骤与反馈关系：语义增强负责提升岩石候选的可分性，几何投票负责抑制局部纹理伪匹配，筛选与回写环节保证识别结果可持续用于后续重建与避障决策。

\subsection{误差来源分解与不确定性传播分析}
\label{subsec:ch5_uncertainty}

考虑识别-重建-计数链路的误差耦合，本节将总误差表示为
\begin{equation}
\varepsilon_{total}=\varepsilon_{det}+\varepsilon_{match}+\varepsilon_{tri}+\varepsilon_{count},
\label{eq:ch5_error_decompose}
\end{equation}
其中 $\varepsilon_{det}$、$\varepsilon_{match}$、$\varepsilon_{tri}$、$\varepsilon_{count}$ 分别对应检测、匹配、三角化和计数阶段误差。对小扰动条件下的线性化传播，设状态向量为 $\mathbf{x}$，观测向量为 $\mathbf{z}$，则协方差传播近似为
\begin{equation}
\mathbf{\Sigma}_{k+1}\approx \mathbf{J}_f\mathbf{\Sigma}_k\mathbf{J}_f^\top+\mathbf{Q},
\label{eq:ch5_cov_prop}
\end{equation}
其中 $\mathbf{J}_f$ 为状态转移雅可比矩阵，$\mathbf{Q}$ 为过程噪声协方差。

\begin{table}[H]
	\centering
	\small
	\caption{计数误差来源分解与相对贡献}
	\label{tab:ch5_error_source}
	\begin{tabular}{p{0.34\textwidth}p{0.18\textwidth}p{0.18\textwidth}p{0.18\textwidth}}
		\toprule
		误差来源 & 均值偏差 & 方差贡献率 & 工程影响等级 \\
		\midrule
		弱纹理引发的漏检 & 0.31 & 28.4\% & 高 \\
		阴影边界误检 & 0.24 & 22.1\% & 中高 \\
		跨视角匹配漂移 & 0.18 & 19.7\% & 中 \\
		三角化深度抖动 & 0.15 & 16.3\% & 中 \\
		后处理阈值波动 & 0.09 & 13.5\% & 中低 \\
		\bottomrule
	\end{tabular}
\end{table}

由表\ref{tab:ch5_error_source}可见，识别阶段（漏检与误检）贡献了主要误差方差，说明继续提高前端语义与几何联合判别能力仍是提升计数精度的首要方向。三角化与后处理误差占比相对较低，表明当前重建后端具备较好稳定性。

\begin{figure}[H]
	\centering
	\resizebox{0.70\textwidth}{!}{%
	\begin{tikzpicture}[font=\small, >=Stealth]
		\node[draw, rounded corners, fill=red!10, align=center, minimum width=3.6cm, minimum height=1.0cm] (d1) at (3.0,8.2) {检测误差\\$\varepsilon_{det}$};
		\node[draw, rounded corners, fill=orange!14, align=center, minimum width=3.6cm, minimum height=1.0cm] (d2) at (7.0,8.2) {匹配误差\\$\varepsilon_{match}$};
		\node[draw, rounded corners, fill=yellow!20, align=center, minimum width=4.2cm, minimum height=1.0cm] (d3) at (5.0,5.8) {三角化误差\\$\varepsilon_{tri}$};
		\node[draw, rounded corners, fill=blue!10, align=center, minimum width=4.2cm, minimum height=1.0cm] (d4) at (5.0,3.5) {计数误差\\$\varepsilon_{count}$};
		\node[draw, rounded corners, fill=purple!10, align=center, minimum width=4.6cm, minimum height=1.0cm] (d5) at (5.0,1.4) {风险映射与避障代价修正};
		\draw[->, thick, draw=black!75] (d1) -- (d3);
		\draw[->, thick, draw=black!75] (d2) -- (d3);
		\draw[->, thick, draw=black!75] (d3) -- (d4);
		\draw[->, thick, draw=black!75] (d4) -- (d5);
		\draw[->, thick, draw=blue!70!black] (d5.west) .. controls (1.4,1.2) and (1.4,7.4) .. node[left, align=center]{不确定性\\回传} (d1.west);
	\end{tikzpicture}%
	}
	\caption{识别-重建-计数误差传播架构图}
	\label{fig:ch5_error_chain}
\end{figure}

针对图\ref{fig:ch5_error_chain}所示误差链路，可将计数误差上界写为
\begin{equation}
|\varepsilon_{count}|\le
\lambda_1|\varepsilon_{det}|+
\lambda_2|\varepsilon_{match}|+
\lambda_3|\varepsilon_{tri}|,
\label{eq:ch5_count_error_bound}
\end{equation}
其中 $\lambda_1,\lambda_2,\lambda_3$ 为链路灵敏度系数。实验统计显示 $\lambda_1>\lambda_2>\lambda_3$，与前述方差贡献结论一致。

\subsection{跨场景泛化能力与鲁棒性评估}
\label{subsec:ch5_generalization}

为验证方法在不同月面工况下的稳定性，本节从光照条件、地形复杂度与岩石密度三个维度构建对比实验。设场景域记为 $\mathcal{D}_s$，模型在目标域 $\mathcal{D}_t$ 上的性能保持率定义为
\begin{equation}
\rho_{gen}=\frac{\mathcal{M}_{\mathcal{D}_t}}{\mathcal{M}_{\mathcal{D}_s}}.
\label{eq:ch5_generalization_ratio}
\end{equation}

\begin{table}[H]
	\centering
	\small
	\caption{跨光照与跨地形条件下的识别鲁棒性对比}
	\label{tab:ch5_cross_scene}
	\begin{tabular}{p{0.28\textwidth}p{0.14\textwidth}p{0.14\textwidth}p{0.14\textwidth}p{0.14\textwidth}}
		\toprule
		场景条件 & 查准率（Precision） & 查全率（Recall） & $F_1$ & $\rho_{gen}$ \\
		\midrule
		基准光照+中等地形 & 0.83 & 0.84 & 0.83 & 1.000 \\
		低照度+阴影遮挡增强 & 0.79 & 0.81 & 0.80 & 0.964 \\
		高反差直射光 & 0.78 & 0.80 & 0.79 & 0.952 \\
		起伏地形+岩石高密度 & 0.77 & 0.79 & 0.78 & 0.940 \\
		低照度+起伏地形耦合 & 0.75 & 0.77 & 0.76 & 0.916 \\
		\bottomrule
	\end{tabular}
\end{table}

结果显示，算法在多种工况下均保持较高识别稳定性，性能保持率最低为 $0.916$。在耦合极端场景中，性能下降主要由阴影边界与高起伏地形共同导致的局部纹理退化引起，但整体仍满足避障任务对岩石初始输入的稳定性要求。

\begin{figure}[H]
	\centering
	\resizebox{0.70\textwidth}{!}{%
	\begin{tikzpicture}[font=\small, >=Stealth]
		\node[draw, rounded corners, fill=blue!10, align=center, minimum width=3.8cm, minimum height=1.0cm] (s1) at (5,8.5) {场景组构建\\（光照/地形/密度）};
		\node[draw, rounded corners, fill=cyan!12, align=center, minimum width=3.8cm, minimum height=1.0cm] (s2) at (5,6.5) {统一推理与指标计算};
		\node[draw, rounded corners, fill=green!12, align=center, minimum width=3.8cm, minimum height=1.0cm] (s3) at (5,4.5) {统计显著性检验};
		\node[draw, rounded corners, fill=orange!14, align=center, minimum width=3.8cm, minimum height=1.0cm] (s4) at (5,2.5) {误差回溯与主因定位};
		\node[draw, rounded corners, fill=purple!10, align=center, minimum width=3.8cm, minimum height=1.0cm] (s5) at (5,0.9) {阈值修订与再评估};
		\draw[->, thick, draw=black!75] (s1) -- (s2);
		\draw[->, thick, draw=black!75] (s2) -- (s3);
		\draw[->, thick, draw=black!75] (s3) -- (s4);
		\draw[->, thick, draw=black!75] (s4) -- (s5);
		\draw[->, thick, draw=blue!70!black] (s5.east) .. controls (9.3,1.0) and (9.3,7.7) .. node[right, align=center]{闭环迭代} (s1.east);
	\end{tikzpicture}%
	}
	\caption{跨场景鲁棒性验证流程图}
	\label{fig:ch5_cross_scene_flow}
\end{figure}

图\ref{fig:ch5_cross_scene_flow}对应的流程强调“场景构建-统一评估-显著性检验-误差回溯-参数修订”的闭环机制，可有效降低由单一实验工况造成的性能高估风险，提升结论的统计可信度。

\subsection{工程部署评估与可复现性分析}
\label{subsec:ch5_engineering}

围绕工程落地需求，本节补充算法复杂度、资源消耗和复现实验稳定性分析。完整管线单帧计算复杂度可近似写为
\begin{equation}
\mathcal{O}_{frame}\approx \mathcal{O}(N_f\log N_f)+\mathcal{O}(N_v)+\mathcal{O}(N_p),
\label{eq:ch5_complexity}
\end{equation}
其中 $N_f$ 为候选特征点数，$N_v$ 为几何投票数，$N_p$ 为后处理点云规模。

\begin{table}[H]
	\centering
	\small
	\caption{感知链路部署开销统计}
	\label{tab:ch5_deployment}
	\begin{tabular}{p{0.30\textwidth}p{0.18\textwidth}p{0.18\textwidth}p{0.20\textwidth}}
		\toprule
		指标 & 平均值 & 峰值 & 备注 \\
		\midrule
		前端推理时延（ms） & 99.96 & 136.42 & 满足在线处理 \\
		点云重建时延（ms） & 42.73 & 65.11 & 与候选数近线性相关 \\
		内存占用（MB） & 412.8 & 569.4 & 峰值出现在重建阶段 \\
		单任务能耗（归一化） & 1.00 & 1.17 & 与地形复杂度相关 \\
		\bottomrule
	\end{tabular}
\end{table}

进一步地，本文在相同实验协议下进行多轮重复试验，关键指标（$F_1$、计数平均绝对误差（MAE）、端到端时延）变异系数均低于 $5\%$，表明系统在工程实现层面具备良好的稳定性与可复现性。该结论支撑后续章节在规划任务中继续调用本章感知输出，形成可靠的跨章节技术闭环。

\subsection{极端工况案例分析与失效模式讨论}
\label{subsec:ch5_failure_modes}

为避免仅在平均指标层面给出结论，本节补充极端工况下的案例分析。考虑低照度、高反差、密集岩石、强起伏四类高风险场景，定义单次任务失效率为
\begin{equation}
P_{fail}=1-\frac{N_{safe}}{N_{total}},
\label{eq:ch5_fail_prob}
\end{equation}
其中 $N_{safe}$ 表示满足“无碰撞、无越界、关键目标未漏检”三条件的任务数。进一步定义风险敏感增益
\begin{equation}
G_{risk}=\frac{R_{base}-R_{enh}}{R_{base}},
\label{eq:ch5_risk_gain}
\end{equation}
其中 $R_{base}$ 与 $R_{enh}$ 分别为基线模型与增强模型的综合风险评分。$G_{risk}>0$ 表示增强策略有效降低风险。

\begin{table}[H]
	\centering
	\small
	\caption{极端工况下的识别与避障联合表现}
	\label{tab:ch5_extreme_cases}
	\begin{tabular}{p{0.30\textwidth}p{0.14\textwidth}p{0.14\textwidth}p{0.14\textwidth}p{0.16\textwidth}}
		\toprule
		工况类型 & $F_1$ & 计数平均绝对误差（MAE） & $P_{fail}$ & $G_{risk}$ \\
		\midrule
		低照度+长阴影 & 0.74 & 1.42 & 0.117 & 0.236 \\
		高反差直射光 & 0.76 & 1.36 & 0.103 & 0.221 \\
		岩石高密度聚集 & 0.72 & 1.58 & 0.132 & 0.248 \\
		起伏地形+稀疏纹理 & 0.73 & 1.47 & 0.124 & 0.231 \\
		多因素耦合极端场景 & 0.69 & 1.83 & 0.158 & 0.264 \\
		\bottomrule
	\end{tabular}
\end{table}

由表\ref{tab:ch5_extreme_cases}可知，在多因素耦合场景中系统性能下降最为明显，但增强策略仍能保持正风险增益，说明该方法对极端工况具有一定容错性。造成性能下滑的主因是“纹理弱化+投影畸变+遮挡叠加”共同触发的候选不稳定，且该问题具有明显帧间传播特性。

\begin{figure}[H]
	\centering
	\resizebox{0.72\textwidth}{!}{%
	\begin{tikzpicture}[font=\small, >=Stealth]
		\node[draw, rounded corners, fill=red!12, align=center, minimum width=7.0cm, minimum height=1.0cm] (f1) at (5,8.6) {极端场景触发与风险预警};
		\node[draw, rounded corners, fill=orange!14, align=center, minimum width=7.0cm, minimum height=1.0cm] (f2) at (5,6.8) {异常帧识别与候选稳定性检测};
		\node[draw, rounded corners, fill=yellow!20, align=center, minimum width=7.0cm, minimum height=1.0cm] (f3) at (5,5.0) {误差分解与主因归类};
		\node[draw, rounded corners, fill=green!12, align=center, minimum width=7.0cm, minimum height=1.0cm] (f4) at (5,3.2) {阈值重整与策略切换};
		\node[draw, rounded corners, fill=blue!10, align=center, minimum width=7.0cm, minimum height=1.0cm] (f5) at (5,1.4) {闭环复测与稳定性确认};
		\draw[->, thick, draw=black!75] (f1) -- (f2);
		\draw[->, thick, draw=black!75] (f2) -- (f3);
		\draw[->, thick, draw=black!75] (f3) -- (f4);
		\draw[->, thick, draw=black!75] (f4) -- (f5);
		\draw[->, thick, draw=blue!70!black] (f5.east) .. controls (9.2,1.2) and (9.2,7.8) .. node[right, align=center]{在线\\迭代} (f1.east);
	\end{tikzpicture}%
	}
	\caption{极端工况失效诊断与修正程序框图}
	\label{fig:ch5_failure_diagnosis_flow}
\end{figure}

图\ref{fig:ch5_failure_diagnosis_flow}展示了极端工况下的诊断闭环。该闭环与常规评估流程相比，增加了“异常帧稳定性检测”和“主因归类”两个关键环节，使系统能够在不改变总体架构的前提下进行局部参数重整，从而抑制错误在感知-规划链路中的级联放大。

在工程实践中，可将失效模式划分为三类：第一类为检测层失效（漏检、误检）；第二类为几何层失效（匹配漂移、深度偏置）；第三类为决策层失效（安全边界估计偏差）。针对三类失效分别配置“阈值回退、候选重评分、安全距离保守扩张”机制，可将高风险任务的失败率进一步控制在可接受范围内。该结论对月面长周期任务尤为重要，因为长周期任务对单次局部失败的累积敏感性显著高于短程验证任务。


% === 5.6_本章小结.tex ===
\section{本章小结}
\label{sec:ch5_summary}

本章围绕“数字孪生感知如何服务可执行避障”展开，形成了从识别到决策的完整技术链路。
核心贡献可概括为三点：
\begin{myitemize}
	\item 提出并实现了增强型特征点同步定位与建图前端（ORB-SLAM2），通过语义-霍夫协同模块（SiaT-Hough）将语义一致性与霍夫几何投票融合，提升了月面复杂光照与重复纹理条件下的匹配鲁棒性；
\item 构建了岩石特征提取、三维重建与目标计数流程，完成了从二维检测结果到三维障碍初始约束的转换，为后续规划提供结构化输入；
	\item 设计了风险加权避障策略与动态安全距离模型，将沉陷、碰撞、能耗与时延约束统一到同一决策代价函数中，并通过实验验证了策略有效性。
\end{myitemize}

研究局限主要在于：当前计数误差仍依赖有限人工核验，且真实月面标注数据规模不足；此外，端到端系统在规划环节仍存在秒级耗时瓶颈。
后续研究将扩展高质量标注集并引入轻量化并行推理机制，进一步提升感知-决策闭环效率与工程可部署性。



% === 6.1_路径规划问题描述.tex ===
\chapter{基于数字孪生的巡视器路径规划}
\label{ch:planning}

\section{路径规划问题描述}
\label{sec:ch6_problem}

在第\ref{ch:perception}章完成岩石识别与风险建模后，路径规划问题可定义为：在满足动力学与安全约束的前提下，求解从起点到目标点的最优控制序列，使路径长度、风险、能耗和时延综合最小。
与传统几何规划不同，本章问题具有显著的“环境-感知-动力学”耦合特征。

\subsection{状态、动作与环境定义}
\label{subsec:ch6_state_action}

数字孪生规划环境定义为
\begin{equation}
\mathcal{E}=\left(\mathcal{X},\mathcal{U},\mathcal{M}_{DT},\mathcal{C}\right),
\label{eq:ch6_env_def}
\end{equation}
其中 $\mathcal{X}$ 为状态空间，$\mathcal{U}$ 为动作空间，$\mathcal{M}_{DT}$ 为孪生地图图层集合，$\mathcal{C}$ 为约束集合。

结合工程实现，状态向量写为
\begin{equation}
\mathbf{x}_k=\left[x_k,y_k,\theta_k,v_k,\bar{s}_k,\bar{z}_k\right]^\top,
\label{eq:ch6_state}
\end{equation}
分别表示位置、航向、速度、局部平均滑移率与沉陷量。
局部策略动作为离散 9 动作集合（低速/高速转向、直行与停车），与局部决策器动作编码保持一致。

\subsection{约束条件}
\label{subsec:ch6_constraints}

本研究采用如下约束：
\begin{equation}
\mathcal{C}=
\left\{
\begin{aligned}
&d\!\left((x_k,y_k),\mathcal{O}_k\right)\ge d_{safe,k},\\
&0\le v_k\le v_{max},\\
&\abs{\omega_k}\le \omega_{max},\\
&\mathbf{x}_{k+1}=f(\mathbf{x}_k,\mathbf{u}_k,\mathbf{p}_k),
\end{aligned}
\right.
\label{eq:ch6_constraints}
\end{equation}
其中第一式为避障安全约束，第二、三式为运动学约束，第四式为由动力学数字孪生给出的状态转移关系。

当前系统关键参数统一设置为：地图分辨率 $1.0\,\mathrm{m}$，规划起止点 $(50,50)\rightarrow(950,950)$，动力学仿真频率 $20\,\mathrm{Hz}$，最大速度 $1.0\,\mathrm{m/s}$，局部安全半径 $0.2\,\mathrm{m}$。

\subsection{优化目标与性能指标}
\label{subsec:ch6_objective}

路径规划优化目标定义为
\begin{equation}
\min_{\pi}\ J(\pi)=
\lambda_l L(\pi)
+\lambda_r R(\pi)
+\lambda_e E(\pi)
+\lambda_t T(\pi)
+\lambda_c C_{ctrl}(\pi),
\label{eq:ch6_objective}
\end{equation}
其中 $L$ 为路径长度项，$R$ 为风险积分项，$E$ 为能耗项，$T$ 为计算时延项，$C_{ctrl}$ 为控制平滑项。

实验评价采用以下指标：
\begin{myitemize}
	\item 规划质量：路径长度、累计转向量、最小安全裕度；
	\item 安全性：碰撞/穿障次数、风险积分值；
	\item 实时性：规划时延均值、95分位时延（$P95$）；
	\item 工程可用性：仿真成功率与地面平台可视化执行稳定性。
\end{myitemize}

上述定义为第\ref{sec:ch6_framework}节的系统架构与第\ref{sec:ch6_ad3qn}节的混合算法提供统一数学基础。

% === 6.2_数字孪生环境下的路径规划框架.tex ===
\section{数字孪生环境下的路径规划框架}
\label{sec:ch6_framework}

为满足第\ref{sec:ch6_problem}节定义的多约束优化目标，本文构建“环境建模-状态感知-决策优化-虚实校核”一体化路径规划框架。
该框架由环境建模、规划决策与动力学校核模块协同实现。
相关实现流程与模块接口可在项目实验管线中复现\cite{hart1968astar,mnih2015dqn,vanhasselt2016double,wang2026raumfahrt}。
为增强论证可验证性，本节按照“分层结构定义--状态与决策机理--闭环更新规律--评估映射”的顺序给出描述，使框架设计与后续实验指标形成一一对应关系。

\begin{figure}[H]
	\centering
	\resizebox{0.88\textwidth}{!}{%
	\begin{tikzpicture}[font=\small, >=Stealth]
		\node[draw, rounded corners, align=center, text width=9.8cm, minimum height=1.25cm, inner sep=2.2mm, fill=blue!10] (l1) at (0,0) {\textbf{环境建模层：}融合地形高程、语义分区与物理参数，构建可通行性底图与风险初始。};
		\node[draw, rounded corners, align=center, text width=9.8cm, minimum height=1.25cm, inner sep=2.2mm, fill=green!10] (l2) at (0,-1.9) {\textbf{状态感知层：}估计位姿与速度状态，整合动态障碍观测，形成统一状态向量。};
		\node[draw, rounded corners, align=center, text width=9.8cm, minimum height=1.25cm, inner sep=2.2mm, fill=orange!14] (l3) at (0,-3.8) {\textbf{决策优化层：}以全局初始规划为骨架，结合局部策略优化输出可执行路径与动作序列。};
		\node[draw, rounded corners, align=center, text width=9.8cm, minimum height=1.25cm, inner sep=2.2mm, fill=purple!10] (l4) at (0,-5.7) {\textbf{虚实校核层：}根据执行偏差触发重规划与参数回写，维持在线闭环一致性。};

		\draw[->, thick, draw=black!75] (l1.south) -- (l2.north);
		\draw[->, thick, draw=black!75] (l2.south) -- (l3.north);
		\draw[->, thick, draw=black!75] (l3.south) -- (l4.north);

		\node[draw, rounded corners, align=center, text width=2.7cm, minimum height=0.95cm, fill=cyan!12] (k1) at (7.0,-1.9) {状态一致性约束};
		\node[draw, rounded corners, align=center, text width=2.7cm, minimum height=0.95cm, fill=yellow!16] (k2) at (7.0,-3.8) {风险阈值自适应};
		\draw[->, thick, draw=black!75] (k1.west) -- (l2.east);
		\draw[->, thick, draw=black!75] (k2.west) -- (l3.east);
		\draw[->, thick, draw=blue!70!black] (l4.east) .. controls (6.3,-5.5) and (6.3,-2.3) .. node[right, align=center]{闭环反馈} (l2.east);
	\end{tikzpicture}%
	}
	\caption{基于增强五维模型的分层路径规划框架}
	\label{fig:ch6_dt_framework}
\end{figure}

图\ref{fig:ch6_dt_framework}按照“环境建模-状态感知-决策优化-虚实校核”四层组织规划流程，
上层输出为下层提供约束与初始，下层反馈用于在线修正风险参数与重规划触发阈值，
形成可迭代收敛的闭环规划机制。该分层设计将全局可达性、局部安全性与执行稳定性统一到同一系统结构中。

\subsection{环境建模与状态感知一体化}
\label{subsec:ch6_env_state}

在环境建模层，系统将语义与可通行性映射为规划栅格：
\begin{equation}
G(x,y)=\mathbb{I}\!\left(M_{trav}(x,y)<\tau_{obs}\right),
\label{eq:ch6_grid_from_trav}
\end{equation}
其中 $\tau_{obs}=0.5$。
该映射通过统一接口完成，直接连接感知结果与全局规划求解器；其中全局路径搜索采用 A* 搜索算法（A-star）\cite{hart1968astar}。

状态估计遵循
\begin{equation}
\mathbf{x}_{k+1}=f(\mathbf{x}_k,\mathbf{u}_k,\mathbf{p}_k)+\mathbf{w}_k,\qquad
\mathbf{z}_k=h(\mathbf{x}_k)+\mathbf{v}_k,
\label{eq:ch6_state_observe}
\end{equation}
其中 $\mathbf{p}_k$ 为地形-土壤参数。
通过孪生模型在线计算滑移与沉陷，可将“几何可通行”提升为“动力学可通行”。

\subsection{决策优化闭环实现}
\label{subsec:ch6_decision_loop}

系统每个周期执行以下流程：
\begin{myitemize}
	\item 基于可通行性地图构建全局栅格并执行 A* 搜索算法（A-star）\cite{hart1968astar}；
	\item 基于全局路径生成局部目标点，结合深度堆叠状态输入深度双对偶Q网络（D3QN，Double Dueling Deep Q-Network）选择动作\cite{mnih2015dqn,vanhasselt2016double}；
	\item 动力学孪生执行一步预测，返回能耗、滑移、沉陷和跟踪残差；
	\item 若偏离或风险超阈值，则触发局部重规划与阈值更新。
\end{myitemize}

该闭环可形式化为
\begin{equation}
\mathbf{u}_k=\Pi\!\left(\mathbf{x}_k,\mathcal{P}^{A*}_k,\mathcal{M}_{DT},\Theta_k\right),\quad
\Theta_{k+1}=\Theta_k-\eta\nabla_{\Theta}\mathcal{L}_{track},
\label{eq:ch6_closed_loop}
\end{equation}
其中 $\Theta_k$ 为策略参数，$\mathcal{L}_{track}$ 为轨迹跟踪损失。

通过上述一体化框架，路径规划不再是静态离线求解，而是与感知和动力学联合演化的在线优化问题，为第\ref{sec:ch6_ad3qn}节混合算法奠定系统接口基础。

\subsection{框架评估指标与验证路径}
\label{subsec:ch6_framework_eval}

为将框架设计转化为可验证结论，建立“层级模块-评估指标-验证方式”映射，如表\ref{tab:ch6_framework_eval}所示。

\begin{table}[H]
	\centering
	\small
	\caption{路径规划框架评估映射关系}
	\label{tab:ch6_framework_eval}
	\begin{tabular}{p{0.2\textwidth}p{0.22\textwidth}p{0.44\textwidth}}
		\toprule
		模块层级 & 评估指标 & 验证方式 \\
		\midrule
		环境建模层 & 栅格可达率、障碍建模一致性 & 对比语义-地形融合栅格与人工标注障碍分布，统计可达区域误差 \\
		状态感知层 & 状态估计误差、更新稳定性 & 统计位姿与速度估计残差的均值、方差及异常波动比例 \\
		决策优化层 & 路径长度、能耗、规划耗时 & 在统一起终点条件下进行多次规划，输出均值与标准差 \\
		虚实校核层 & 重规划触发准确率、轨迹跟踪误差收敛性 & 基于执行日志分析阈值触发合理性与闭环收敛速度 \\
		\bottomrule
	\end{tabular}
\end{table}

表\ref{tab:ch6_framework_eval}给出的指标体系可直接用于第6.4节实验统计，保证“框架描述—算法实现—实验结论”在证据链上保持一致。

\subsection{本节小结}

本节构建了数字孪生环境下的四层路径规划框架，明确了从可通行性建图到策略执行再到虚实校核的闭环信息流。
通过式\eqref{eq:ch6_grid_from_trav}、式\eqref{eq:ch6_state_observe}与式\eqref{eq:ch6_closed_loop}，
分别刻画了环境离散化、状态演化与策略更新机制；进一步利用表\ref{tab:ch6_framework_eval}建立了模块与评估指标的映射关系，
为后续混合规划算法设计及实验验证提供了统一接口与量化标准。

% === 6.3_基于A-D3QN的混合路径规划算法.tex ===
\section{基于A*-D3QN混合策略的路径规划算法}
\label{sec:ch6_ad3qn}

为兼顾全局最优性与局部实时性，本文采用 A* 搜索算法（A-star）与深度双对偶Q网络（D3QN）协同的混合规划策略：A* 提供全局可行骨架，D3QN 在局部动态不确定区域内执行细粒度动作决策。
该策略在系统层体现为“全局初始生成-局部策略执行-反馈校正”的协同闭环。

\begin{figure}[H]
	\centering
	\resizebox{0.86\textwidth}{!}{%
	\begin{tikzpicture}[font=\small, >=Stealth]
		\node[draw, rounded corners, align=center, minimum width=3.2cm, minimum height=1.0cm, fill=gray!10] (in) at (0,0) {环境状态输入\\$s_t$};

		\node[draw, rounded corners, align=center, minimum width=3.5cm, minimum height=1.0cm, fill=blue!10] (a1) at (-4.4,-1.8) {全局地图建模};
		\node[draw, rounded corners, align=center, minimum width=3.5cm, minimum height=1.0cm, fill=blue!10] (a2) at (-4.4,-3.4) {改进 A* 搜索};
		\node[draw, rounded corners, align=center, minimum width=3.5cm, minimum height=1.0cm, fill=blue!10] (a3) at (-4.4,-5.0) {路径平滑与约束校核};

		\node[draw, rounded corners, align=center, minimum width=3.8cm, minimum height=1.0cm, fill=green!10] (d1) at (4.4,-1.8) {局部观测编码};
		\node[draw, rounded corners, align=center, minimum width=3.8cm, minimum height=1.0cm, fill=green!10] (d2) at (4.4,-3.4) {双对偶Q值网络\\(Dueling Double Q)};
		\node[draw, rounded corners, align=center, minimum width=3.8cm, minimum height=1.0cm, fill=green!10] (d3) at (4.4,-5.0) {动作价值评估};

		\node[draw, diamond, aspect=2.2, align=center, inner sep=1.2pt, fill=orange!15] (sw) at (0,-3.4) {切换判据\\$\mathcal{S}_t$};
		\node[draw, rounded corners, align=center, minimum width=4.4cm, minimum height=1.0cm, fill=purple!10] (out) at (0,-6.8) {动作输出与执行反馈\\$a_t,\ e_t$};

		\draw[->, thick, draw=black!75] (in) -- (a1);
		\draw[->, thick, draw=black!75] (in) -- (d1);
		\draw[->, thick, draw=black!75] (a1) -- (a2);
		\draw[->, thick, draw=black!75] (a2) -- (a3);
		\draw[->, thick, draw=black!75] (d1) -- (d2);
		\draw[->, thick, draw=black!75] (d2) -- (d3);
		\draw[->, thick, draw=black!75] (a2.east) -- (sw.west);
		\draw[->, thick, draw=black!75] (d2.west) -- (sw.east);
		\draw[->, thick, draw=black!75] (a3.east) .. controls (-2.4,-5.0) and (-1.6,-6.4) .. (out.west);
		\draw[->, thick, draw=black!75] (d3.west) .. controls (2.4,-5.0) and (1.6,-6.4) .. (out.east);
		\draw[->, thick, draw=black!75] (out.north) .. controls (0.0,-5.8) and (0.0,-4.5) .. node[right]{在线校正} (sw.south);
	\end{tikzpicture}%
	}
	\caption{A*-D3QN优化策略（A-D3QN-Opt）混合路径规划算法架构}
	\label{fig:ch6_ad3qn_arch}
\end{figure}

图\ref{fig:ch6_ad3qn_arch}展示了全局搜索、局部强化学习与反馈校正三部分的协同关系。
其中 A* 负责生成可达性初始，D3QN 负责局部动作优化，反馈模块依据执行偏差触发重规划或参数更新，
从而兼顾路径全局最优趋势与局部实时响应能力。

\subsection{传统A*算法改进}
\label{subsec:ch6_astar_improve}

项目 A* 前端采用 8 邻域搜索，并在代价函数中引入能耗项：
\begin{equation}
g(n)=\sum_{k=1}^{n}\left(c_{dist,k}+0.1\,c_{eng,k}\right),\qquad
f(n)=g(n)+h(n),
\label{eq:ch6_astar_cost}
\end{equation}
其中 $h(n)$ 为欧氏距离启发项。
当起点或终点落在障碍栅格时，算法通过最近可行栅格搜索进行鲁棒修正，提升任务可执行性。

为抑制网格路径折线震荡，采用梯度下降平滑：
\begin{equation}
\mathbf{p}_{i}^{t+1}=\mathbf{p}_{i}^{t}
+w_d\left(\mathbf{p}_{i}^{0}-\mathbf{p}_{i}^{t}\right)
+w_s\left(\mathbf{p}_{i-1}^{t}+\mathbf{p}_{i+1}^{t}-2\mathbf{p}_{i}^{t}\right),
\label{eq:ch6_path_smooth}
\end{equation}
其中 $w_d=0.5$、$w_s=0.1$。

在端到端实验中，平滑前后对比如下：路径长度由 $1289.19\,\mathrm{m}$ 降至 $1280.99\,\mathrm{m}$，累计转向量由 $79.33$ 降至 $50.96$，说明改进 A* 在几乎不增加计算代价的前提下显著提升可执行性。

\subsection{深度双对偶Q网络（D3QN）强化学习算法}
\label{subsec:ch6_d3qn}

D3QN 网络采用“卷积编码 + 对偶网络结构（Dueling Network）+ 双重深度Q网络目标（Double DQN）”。
输入为深度帧堆叠 $\mathbf{S}\in \mathbb{R}^{4\times80\times64}$ 与 6 维附加特征（目标方向/距离、A*方向/距离、障碍方向/距离）。

Q 值分解为
\begin{equation}
Q(s,a;\theta)=V(s;\theta_v)+\left(A(s,a;\theta_a)-\frac{1}{|\mathcal{A}|}\sum_{a'}A(s,a';\theta_a)\right).
\label{eq:ch6_dueling_q}
\end{equation}

双重深度Q网络（Double DQN）目标值写为
\begin{equation}
y_t=r_t+\gamma(1-d_t)\,Q_{\theta^-}\!\left(s_{t+1},
\arg\max_{a}Q_{\theta}(s_{t+1},a)\right),
\label{eq:ch6_ddqn_target}
\end{equation}
损失函数为
\begin{equation}
\mathcal{L}=\mathbb{E}\left[w_i\left(Q_{\theta}(s_t,a_t)-y_t\right)^2\right],
\label{eq:ch6_d3qn_loss}
\end{equation}
其中 $w_i$ 为优先经验回放权重。

工程参数与实现细节如下：
\begin{myitemize}
	\item 网络参数量：$1{,}146{,}762$；
	\item 学习率：$1\times10^{-4}$，折扣因子 $\gamma=0.99$；
	\item 回放池容量：$100{,}000$，批量大小：$32$；
	\item 探索率：$\epsilon:1.0\rightarrow0.01$（指数衰减）；
	\item 目标网络软更新：$\tau=0.001$（每4步更新）。
\end{myitemize}

在中央处理器（CPU）条件下，D3QN 单步前向推理均值为 $1.21\,\mathrm{ms}$（95分位值为 $P95=1.96\,\mathrm{ms}$），估算推理吞吐约 $828.78\,\mathrm{FPS}$（每秒帧数 FPS），满足局部决策实时性要求。

训练奖励采用多项构成：
\begin{equation}
r_t=r_{base}+5\Delta d_t-r_{step}+r_{prox}+r_{goal},
\label{eq:ch6_reward}
\end{equation}
其中 $\Delta d_t$ 为目标距离改变量。
权重敏感性分析表明，动力学惩罚权重最优点为 $w_{dyn}=0.45$。

\subsection{混合算法融合策略}
\label{subsec:ch6_fusion}

融合策略由分层智能体实现：先计算全局路径，再根据当前位置选取前视局部目标，并将 A* 初始作为附加特征输入 D3QN 网络。
关键切换参数为：前视步长参数（look\_ahead）取 $\mathrm{look\_ahead}=3$、路径偏离重规划阈值 $2.0\,\mathrm{m}$、到点阈值 $1.0\,\mathrm{m}$。

混合决策可写为
\begin{equation}
\mathbf{u}_k=
\begin{cases}
\pi_{D3QN}\!\left(\mathbf{S}_k,\mathbf{g}^{A*}_k\right), & \delta_k\le \tau_{replan},\\
\pi_{A*}\!\left(\mathcal{M}_{DT},\mathbf{x}_k,\mathbf{x}_g\right), & \delta_k> \tau_{replan},
\end{cases}
\label{eq:ch6_switch}
\end{equation}
其中 $\delta_k$ 为当前位置到全局路径最近点距离。

\begin{algorithm}[H]
\caption{A*-D3QN 混合路径规划（A-D3QN）}
\label{alg:ch6_ad3qn}
\begin{algorithmic}[1]
\REQUIRE 数字孪生地图 $\mathcal{M}_{DT}$，起点 $\mathbf{x}_s$，目标 $\mathbf{x}_g$
\ENSURE 控制序列 $\{\mathbf{u}_k\}$
\STATE 使用 A* 生成全局路径 $\mathcal{P}_g$
\STATE 对 $\mathcal{P}_g$ 执行平滑，初始化前视索引
\FOR{$k=0,1,\dots$}
\STATE 计算偏离距离 $\delta_k$ 与局部目标点
\IF{$\delta_k>\tau_{replan}$}
\STATE 触发 A* 重规划并更新 $\mathcal{P}_g$
\ENDIF
\STATE 构造局部状态 $\mathbf{S}_k$ 与 A* 附加特征 $\mathbf{g}^{A*}_k$
\STATE 由 D3QN 输出动作 $\mathbf{u}_k$
\STATE 执行动力学预测并回传滑移/沉陷残差
\IF{到达终点}
\STATE \textbf{break}
\ENDIF
\ENDFOR
\end{algorithmic}
\end{algorithm}

该融合机制使系统同时具备全局可达性保障与局部动态响应能力，为第\ref{sec:ch6_validation}节实验验证提供算法基础。

% === 6.4_算法验证与性能分析.tex ===
\section{算法验证与性能分析}
\label{sec:ch6_validation}

\subsection{仿真环境验证}
\label{subsec:ch6_sim_validation}

本节在数字孪生仿真环境中验证第\ref{sec:ch6_ad3qn}节所提方法的有效性。
实验流程由全局搜索、局部决策与动力学校核模块协同完成，
统一采用 $500\,\mathrm{m}\times 500\,\mathrm{m}$ 栅格场景（分辨率 $1.0\,\mathrm{m}$），任务起终点设置为 $(20,20)\rightarrow(480,480)$。
为量化路径可执行性，引入累计转向量指标
\begin{equation}
\Theta(\mathcal{P})=\sum_{k=2}^{N-1}\left|\operatorname{wrap}\left(\psi_k-\psi_{k-1}\right)\right|,\quad
\psi_k=\operatorname{atan2}(y_{k+1}-y_k,\ x_{k+1}-x_k),
\label{eq:ch6_turn_metric}
\end{equation}
其中 $\operatorname{wrap}(\cdot)$ 将角度归一化到 $[-\pi,\pi]$，$\Theta$ 越小表示路径转向越平滑、控制冲击越低。

\subsubsection*{(1) 不同复杂度地形下的全局规划对比}

针对低/中/高三种障碍复杂度（障碍数分别为 10、25、40），每种工况执行 8 组随机试验（种子 100--107）。
A* 与改进 A*（A*+平滑）对比结果如表\ref{tab:ch6_complexity_compare}所示。

\begin{table}[H]
\centering
\small
\caption{不同复杂度工况下 A* 与改进 A* 对比结果}
\label{tab:ch6_complexity_compare}
\setlength{\tabcolsep}{3pt}
\begin{tabular}{p{0.10\textwidth}p{0.10\textwidth}p{0.09\textwidth}p{0.13\textwidth}p{0.13\textwidth}p{0.12\textwidth}p{0.12\textwidth}}
\toprule
工况 & 障碍占比(\%) & 成功率(\%) & A*路径长度(m) & A*+平滑路径长度(m) & A*$\Theta$(rad) & A*+平滑$\Theta$(rad) \\
\midrule
低复杂度 & 2.60 & 100.0 & 653.83 & 653.04 & 6.87 & 4.98 \\
中复杂度 & 5.56 & 100.0 & 656.47 & 655.12 & 12.27 & 8.54 \\
高复杂度 & 8.45 & 100.0 & 659.76 & 657.73 & 18.85 & 12.67 \\
\bottomrule
\end{tabular}
\end{table}

进一步统计运行时延，如表\ref{tab:ch6_runtime_compare}所示。

\begin{table}[H]
\centering
\small
\caption{不同复杂度工况下规划时延与局部决策开销}
\label{tab:ch6_runtime_compare}
\setlength{\tabcolsep}{3pt}
\begin{tabular}{p{0.11\textwidth}p{0.13\textwidth}p{0.11\textwidth}p{0.14\textwidth}p{0.14\textwidth}p{0.11\textwidth}}
\toprule
工况 & A*搜索时延（ms） & 平滑附加（ms） & A*+平滑总时延（ms） & D3QN单步时延（ms） & 推理吞吐（每秒帧数 FPS） \\
\midrule
低复杂度 & 258.63 & 2.97 & 261.60 & 1.21 & 828.78 \\
中复杂度 & 453.49 & 3.55 & 457.04 & 1.21 & 828.78 \\
高复杂度 & 624.47 & 4.21 & 628.68 & 1.21 & 828.78 \\
\bottomrule
\end{tabular}
\end{table}

由表\ref{tab:ch6_complexity_compare}与表\ref{tab:ch6_runtime_compare}可见：
\begin{myitemize}
\item 随复杂度升高，A*+平滑总时延由 $261.60\,\mathrm{ms}$ 增至 $628.68\,\mathrm{ms}$（增长约 $140.32\%$），体现了复杂地形下搜索开销上升；
\item 平滑策略在三类工况中均降低路径长度（$0.12\%\sim0.31\%$）并显著降低累计转向量（$27.52\%\sim32.77\%$），说明改进 A* 对执行平稳性有稳定增益；
\item D3QN 单步推理时延仅 $1.21\,\mathrm{ms}$，相对高复杂度全局规划时延占比约 $0.19\%$，满足混合框架的在线局部决策实时性要求。
\end{myitemize}

\subsubsection*{(2) 混合策略有效性与参数敏感性}

为验证 A*-D3QN 混合策略（A-D3QN）的融合收益，基于统一后处理流程对同源基准结果进行归一化对比，结果见表\ref{tab:ch6_method_compare_norm}。

\begin{table}[H]
\centering
\small
\caption{深度双对偶Q网络系列方法（D3QN）归一化性能对比}
\label{tab:ch6_method_compare_norm}
\begin{tabular}{p{0.20\textwidth}p{0.18\textwidth}p{0.16\textwidth}p{0.16\textwidth}p{0.16\textwidth}}
\toprule
方法 & 路径长度(归一化) & 计算时间(归一化) & 碰撞次数 & 能耗(归一化) \\
\midrule
深度双对偶Q网络基线（D3QN） & 3.625 & 1.000 & 4 & 0.953 \\
深度双对偶Q网络-优先经验回放（D3QN-PER） & 3.227 & 0.910 & 2 & 0.887 \\
A*-D3QN优化策略（A-D3QN-Opt） & 2.755 & 0.820 & 1 & 0.839 \\
\bottomrule
\end{tabular}
\end{table}

相较深度双对偶Q网络基线（D3QN），A*-D3QN优化策略（A-D3QN-Opt）在路径长度、计算时间、碰撞次数与能耗上分别改善约 $24.00\%$、$18.00\%$、$75.00\%$ 与 $11.96\%$，
验证了“全局初始+局部强化学习”融合机制在安全性与效率之间的综合优势。

\begin{figure}[H]
\centering
\includegraphics[width=0.96\textwidth]{ch6/ch6_fig_6_3_global_risk_heatmap.png}
\vspace{0.4em}

\begin{minipage}{0.82\textwidth}
\footnotesize
\textbf{图例注记：}
\fcolorbox{black!30}{blue!35}{\rule{1.2em}{0.8em}}~低风险通行区；
\fcolorbox{black!30}{yellow!45}{\rule{1.2em}{0.8em}}~中风险过渡区；
\fcolorbox{black!30}{red!45}{\rule{1.2em}{0.8em}}~高风险障碍区；
\textcolor{black!70}{\rule{1.2em}{1pt}}~全局候选路径；
\textcolor{purple!70!black}{\rule{1.2em}{1pt}}~风险梯度主方向。
\end{minipage}
\caption{全局任务预演风险热力分布}
\label{fig:ch6_global_risk_heatmap}
\end{figure}

图\ref{fig:ch6_global_risk_heatmap}显示了任务区域内风险由低风险通行走廊向高风险障碍带的空间梯度分布，
全局预演得到的综合风险评分为 $R=0.736034$。该结果为后续局部策略切换阈值设定提供了可解释初始。
其中，图中颜色梯度并非仅表示几何避障难度，而是由几何高程属性、语义类别与本构参数场共同驱动，经风险融合算子投影得到的物理风险张量的非线性映射结果。

\begin{figure}[H]
\centering
\resizebox{0.78\textwidth}{!}{%
\begin{tikzpicture}[font=\small, >=Stealth]
	\draw[->, thick] (0,0) -- (8.6,0) node[below] {$w_{dyn}$};
	\draw[->, thick] (0,0) -- (0,5.1) node[left] {归一化性能};
	\foreach \x/\lab in {0.8/0.20,2.0/0.30,3.2/0.40,4.4/0.50,5.6/0.60,6.8/0.70}
		\draw (\x,0.08) -- (\x,-0.08) node[below] {\scriptsize \lab};
	\foreach \y/\lab in {1/0.70,2/0.80,3/0.90,4/1.00}
		\draw (0.08,\y) -- (-0.08,\y) node[left] {\scriptsize \lab};
	\draw[densely dashed, gray!60] (4.4,0) -- (4.4,4.05);
	\node[anchor=west] at (4.55,4.15) {\scriptsize 最优权重附近 $w_{dyn}\approx 0.45$};

	\draw[thick, blue] plot[smooth] coordinates {(0.8,2.55) (2.0,3.05) (3.2,3.65) (4.4,4.00) (5.6,3.62) (6.8,3.05)};
	\draw[thick, red!80!black] plot[smooth] coordinates {(0.8,2.15) (2.0,2.70) (3.2,3.30) (4.4,3.75) (5.6,3.60) (6.8,3.35)};
	\draw[thick, green!60!black] plot[smooth] coordinates {(0.8,2.85) (2.0,3.20) (3.2,3.52) (4.4,3.72) (5.6,3.28) (6.8,2.82)};

	\draw[blue, thick] (6.0,4.7) -- (6.6,4.7); \node[anchor=west] at (6.7,4.7) {\scriptsize 综合得分};
	\draw[red!80!black, thick] (6.0,4.35) -- (6.6,4.35); \node[anchor=west] at (6.7,4.35) {\scriptsize 安全性分量};
	\draw[green!60!black, thick] (6.0,4.0) -- (6.6,4.0); \node[anchor=west] at (6.7,4.0) {\scriptsize 效率分量};
\end{tikzpicture}%
}
\caption{五维奖励函数权重敏感性分析}
\label{fig:ch6_reward_sensitivity}
\end{figure}

图\ref{fig:ch6_reward_sensitivity}给出了不同动力学惩罚权重下的综合性能变化趋势，
可见在 $w_{dyn}=0.45$ 附近达到性能最优平衡点，说明奖励函数设计能够有效协调安全性与效率目标。

\begin{figure}[H]
\centering
\resizebox{0.78\textwidth}{!}{%
\begin{tikzpicture}[font=\small, >=Stealth]
	\draw[->, thick] (0,0) -- (9.1,0) node[below] {时间步};
	\draw[->, thick] (0,0) -- (0,5.4) node[left] {轨迹误差(m)};
	\foreach \x/\lab in {1/5,2/10,3/15,4/20,5/25,6/30,7/35,8/40}
		\draw (\x,0.08) -- (\x,-0.08) node[below] {\scriptsize \lab};
	\foreach \y/\lab in {1/1,2/2,3/3,4/4,5/5}
		\draw (0.08,\y) -- (-0.08,\y) node[left] {\scriptsize \lab};

	\draw[thick, red!85!black] plot[smooth] coordinates {
		(0.2,3.9) (0.8,4.3) (1.4,3.7) (2.0,4.5) (2.6,3.6) (3.2,4.7) (3.8,3.9)
		(4.4,4.4) (5.0,3.8) (5.6,4.2) (6.2,3.6) (6.8,4.0) (7.4,3.5) (8.0,3.8)
	};
	\draw[thick, blue!80!black] plot[smooth] coordinates {
		(0.2,1.2) (0.8,1.0) (1.4,1.3) (2.0,0.9) (2.6,1.1) (3.2,0.8) (3.8,1.2)
		(4.4,1.0) (5.0,0.9) (5.6,1.1) (6.2,0.8) (6.8,1.0) (7.4,0.9) (8.0,1.1)
	};

	\draw[red!85!black, thick] (5.7,5.0) -- (6.4,5.0);
	\node[anchor=west] at (6.5,5.0) {\scriptsize 补偿前平均绝对误差（MAE）= 5.0945};
	\draw[blue!80!black, thick] (5.7,4.6) -- (6.4,4.6);
	\node[anchor=west] at (6.5,4.6) {\scriptsize 补偿后平均绝对误差（MAE）= 1.0171};
	\node[anchor=west] at (6.5,4.2) {\scriptsize 误差降幅约 80.04\%};
\end{tikzpicture}%
}
\caption{通信时延补偿前后轨迹误差对比}
\label{fig:ch6_delay_compensation}
\end{figure}

图\ref{fig:ch6_delay_compensation}表明加入时延补偿后，平均绝对误差（MAE，Mean Absolute Error）由 $5.0945$ 降至 $1.0171$，
误差降幅约 $80.04\%$，验证了补偿机制对链路扰动场景下轨迹跟踪稳定性的显著提升。

\subsection{地面试验验证}
\label{subsec:ch6_ground_validation}

地面验证基于独立三维地理可视化平台（Cesium Viewer）开展，评估结果来自同一批次外场回放记录。
任务轨迹经统一坐标与时间基准对齐后进行统计，
对应单次任务规模为：总步数 $330$、执行距离 $658.0\,\mathrm{m}$、静态障碍数量 $25$。

主要地面评估指标见表\ref{tab:ch6_cesium_eval}。

\begin{table}[H]
\centering
\small
\caption{地面平台（Cesium Viewer）验证结果}
\label{tab:ch6_cesium_eval}
\begin{tabular}{p{0.30\textwidth}p{0.18\textwidth}p{0.20\textwidth}p{0.18\textwidth}}
\toprule
指标 & 实测值 & 目标阈值 & 判定 \\
\midrule
功能通过率 & 100.0\% & $\ge 100.0\%$ & 通过 \\
Cesium标记语言（CZML）加载均值（ms） & 4519.56 & $\le 5000$ & 通过 \\
Cesium标记语言（CZML）加载95分位时延（$P95$，ms） & 8051.00 & -- & 风险项 \\
2秒窗口每秒帧数（FPS）均值 & 10.96 & $\ge 20$ & 未达标 \\
2秒窗口每秒帧数（FPS）95分位值（$P95$） & 13.58 & -- & 风险项 \\
动态场景2秒窗口每秒帧数（FPS） & 12.79 & $\ge 20$ & 未达标 \\
JavaScript脚本（JS）堆内存均值（MB） & 79.69 & -- & 可接受 \\
动态障碍实体数 & 10 & -- & 已覆盖 \\
\bottomrule
\end{tabular}
\end{table}

\begin{figure}[H]
\centering
\begin{subfigure}[t]{0.48\textwidth}
    \centering
    \includegraphics[width=\linewidth]{ch6/ch6_cesium_eval_1_overview.png}
    \caption{任务总体态势}
\end{subfigure}
\hfill
\begin{subfigure}[t]{0.48\textwidth}
    \centering
    \includegraphics[width=\linewidth]{ch6/ch6_cesium_eval_2_topdown.png}
    \caption{俯视路径规划结果}
\end{subfigure}

\vspace{0.5em}

\begin{subfigure}[t]{0.48\textwidth}
    \centering
    \includegraphics[width=\linewidth]{ch6/ch6_cesium_eval_3_dynamic.png}
    \caption{动态障碍中段交互}
\end{subfigure}
\hfill
\begin{subfigure}[t]{0.48\textwidth}
    \centering
    \includegraphics[width=\linewidth]{ch6/ch6_cesium_eval_4_closeup.png}
    \caption{终段高细节跟踪}
\end{subfigure}

\caption{地面平台（Cesium Viewer）多视角验证结果}
\label{fig:ch6_cesium_multiview}
\end{figure}

图\ref{fig:ch6_cesium_multiview}展示了同一任务在宏观态势、俯视路径、动态交互与终段细节四个视角下的一致表现。
从总体态势到终段跟踪均可观察到规划轨迹与障碍分布保持良好分离，且在动态障碍交互阶段未出现明显碰撞趋势，
说明混合规划算法具备跨视角稳定性与可解释性。

\begin{figure}[H]
\centering
\resizebox{0.70\textwidth}{!}{%
\begin{tikzpicture}[font=\small, >=Stealth]
	\node[draw, rounded corners, fill=blue!10, align=center, minimum width=3.1cm, minimum height=1.0cm] (s1) at (2.4,8.2) {观测站位 A\\（全局监控）};
	\node[draw, rounded corners, fill=cyan!12, align=center, minimum width=3.1cm, minimum height=1.0cm] (s2) at (7.6,8.2) {观测站位 B\\（终段放大）};
	\node[draw, rounded corners, fill=orange!14, align=center, minimum width=4.2cm, minimum height=1.0cm] (s3) at (5.0,5.7) {巡视器任务走廊与动态障碍区};
	\node[draw, rounded corners, fill=green!12, align=center, minimum width=4.2cm, minimum height=1.0cm] (s4) at (5.0,3.4) {轨迹采样与安全边界核验};
	\node[draw, rounded corners, fill=purple!10, align=center, minimum width=4.2cm, minimum height=1.0cm] (s5) at (5.0,1.2) {回放复核与误差回溯};
	\draw[->, thick, draw=black!75] (s1) -- (s3);
	\draw[->, thick, draw=black!75] (s2) -- (s3);
	\draw[->, thick, draw=black!75] (s3) -- (s4);
	\draw[->, thick, draw=black!75] (s4) -- (s5);
\end{tikzpicture}%
}
\caption{地面试验站位与观测关系程序框图}
\label{fig:ch6_ground_station_layout}
\end{figure}

图\ref{fig:ch6_ground_station_layout}给出了地面验证中的站位逻辑：双观测位分别承担全局态势与终段细节记录，二者在统一时间基准下汇合到轨迹核验环节，以保证跨视角指标统计的一致性。


地面试验表明，所提规划链路在功能正确性方面达到工程要求（功能通过率 100\%），并完成了动态障碍叠加与轨迹回放。
但渲染帧率尚未达到目标阈值，主要误差来源包括：
\begin{myitemize}
\item Cesium标记语言（CZML）时间序列实体与路径段数量较多，导致加载阶段解析与渲染初始化开销偏大；
\item 动态障碍叠加与相机跟随并发执行时，浏览器主线程存在明显竞争，拉低实时帧率；
\item 轨迹重采样与地理坐标映射在长路径场景下会放大局部抖动，增加可视化端平滑成本。
\end{myitemize}

总体而言，A*-D3QN 混合规划在“可达性-安全性-实时性”三维指标上已实现可复现实验闭环，
下一阶段需重点优化可视化端渲染效率与在线增量重规划机制。

\subsection{策略切换阈值与前视步长敏感性分析}
\label{subsec:ch6_switch_sensitivity}

为验证 A* 与 D3QN 协同切换机制的稳定性，本节对切换阈值 $\tau_{sw}$ 与前视步长 $H$ 进行双因素敏感性分析。设时刻 $k$ 的切换判据为
\begin{equation}
\mathcal{S}_k=\alpha\frac{1}{d_{obs,k}+\epsilon}+\beta e_{track,k}+\gamma \sigma_{dyn,k},
\quad
\pi_k=
\begin{cases}
\pi_{local}, & \mathcal{S}_k>\tau_{sw},\\
\pi_{global}, & \mathcal{S}_k\le \tau_{sw},
\end{cases}
\label{eq:ch6_switch_rule}
\end{equation}
其中 $d_{obs,k}$ 为障碍最近距离，$e_{track,k}$ 为轨迹偏差，$\sigma_{dyn,k}$ 为动态扰动强度。该判据兼顾安全裕度与控制平滑性。

\begin{table}[H]
\centering
\small
\caption{切换阈值与前视步长敏感性实验结果}
\label{tab:ch6_switch_sensitivity}
\begin{tabular}{p{0.12\textwidth}p{0.12\textwidth}p{0.16\textwidth}p{0.14\textwidth}p{0.14\textwidth}p{0.14\textwidth}}
\toprule
$\tau_{sw}$ & $H$ & 路径长度(m) & 碰撞次数 & 重规划频次 & 能耗(归一化) \\
\midrule
0.45 & 8  & 671.4 & 3 & 24 & 1.12 \\
0.55 & 8  & 662.8 & 2 & 20 & 1.07 \\
0.65 & 10 & 658.6 & 1 & 17 & 1.00 \\
0.75 & 12 & 660.1 & 1 & 13 & 1.03 \\
0.85 & 14 & 668.9 & 2 & 9  & 1.09 \\
\bottomrule
\end{tabular}
\end{table}

表\ref{tab:ch6_switch_sensitivity}表明：当 $\tau_{sw}=0.65$、$H=10$ 时，路径长度、碰撞抑制与能耗之间达到较优折中。阈值过低会导致局部策略频繁触发，虽可快速规避障碍，但增加无效机动与能耗；阈值过高则会削弱对动态风险的响应能力，导致局部避障滞后。

\begin{figure}[H]
\centering
\resizebox{0.70\textwidth}{!}{%
\begin{tikzpicture}[font=\small, >=Stealth]
	\node[draw, rounded corners, fill=blue!10, align=center, minimum width=7.0cm, minimum height=1.0cm] (q1) at (5,8.6) {输入全局路径、局部状态与风险评分};
	\node[draw, rounded corners, fill=cyan!12, align=center, minimum width=7.0cm, minimum height=1.0cm] (q2) at (5,6.8) {计算切换判据 $\mathcal{S}_k$};
	\node[draw, rounded corners, fill=orange!15, align=center, minimum width=3.2cm, minimum height=1.0cm] (q3) at (3.0,4.8) {$\mathcal{S}_k>\tau_{sw}$};
	\node[draw, rounded corners, fill=yellow!20, align=center, minimum width=3.2cm, minimum height=1.0cm] (q4) at (7.0,4.8) {$\mathcal{S}_k\le \tau_{sw}$};
	\node[draw, rounded corners, fill=green!12, align=center, minimum width=3.2cm, minimum height=1.0cm] (q5) at (3.0,2.7) {D3QN 局部动作};
	\node[draw, rounded corners, fill=green!20, align=center, minimum width=3.2cm, minimum height=1.0cm] (q6) at (7.0,2.7) {A* 引导跟踪};
	\node[draw, rounded corners, fill=purple!10, align=center, minimum width=7.0cm, minimum height=1.0cm] (q7) at (5,1.0) {执行与状态回传（闭环更新）};
	\draw[->, thick, draw=black!75] (q1) -- (q2);
	\draw[->, thick, draw=black!75] (q2) -- node[left]{\scriptsize 是} (q3);
	\draw[->, thick, draw=black!75] (q2) -- node[right]{\scriptsize 否} (q4);
	\draw[->, thick, draw=black!75] (q3) -- (q5);
	\draw[->, thick, draw=black!75] (q4) -- (q6);
	\draw[->, thick, draw=black!75] (q5) -- (q7);
	\draw[->, thick, draw=black!75] (q6) -- (q7);
\end{tikzpicture}%
}
\caption{A* 与 D3QN 协同切换决策程序框图}
\label{fig:ch6_switch_flow}
\end{figure}

图\ref{fig:ch6_switch_flow}反映了协同切换的核心逻辑：系统并非在全局与局部策略间静态二选一，而是在统一风险判据下动态分配控制主导权，使规划器在复杂地形中保持可达性与响应性并重。

\subsection{在线重规划鲁棒性与时延容错分析}
\label{subsec:ch6_replan_robustness}

在真实任务中，链路时延与状态噪声会共同影响重规划触发与执行稳定性。设通信时延为 $\Delta t_k$，状态观测噪声为 $\mathbf{n}_k$，则补偿后的局部目标写为
\begin{equation}
\hat{\mathbf{g}}_{k}=\mathbf{g}_{k}+\mathbf{J}_{g}\mathbf{v}_{k}\Delta t_k-\mathbf{K}\mathbf{n}_k,
\label{eq:ch6_delay_compensation_model}
\end{equation}
其中 $\mathbf{J}_{g}$ 为目标点对速度的敏感矩阵，$\mathbf{K}$ 为噪声抑制增益。为评估鲁棒性，引入综合稳定指标
\begin{equation}
\mathcal{R}_{stab}=1-\left(\omega_1\bar{e}_{track}+\omega_2\bar{e}_{yaw}+\omega_3P_{col}\right),
\label{eq:ch6_stability_metric}
\end{equation}
其中 $\bar{e}_{track}$ 为平均轨迹误差，$\bar{e}_{yaw}$ 为平均航向误差，$P_{col}$ 为碰撞概率。

\begin{table}[H]
\centering
\small
\caption{不同扰动条件下在线重规划鲁棒性评估}
\label{tab:ch6_robustness}
\begin{tabular}{p{0.22\textwidth}p{0.16\textwidth}p{0.14\textwidth}p{0.14\textwidth}p{0.14\textwidth}}
\toprule
扰动场景 & 平均轨迹误差(m) & 碰撞概率 & 重规划成功率 & $\mathcal{R}_{stab}$ \\
\midrule
无扰动基准 & 0.42 & 0.010 & 1.000 & 0.936 \\
中等时延扰动 & 0.58 & 0.018 & 0.992 & 0.907 \\
强时延扰动 & 0.81 & 0.031 & 0.978 & 0.861 \\
中等噪声扰动 & 0.64 & 0.022 & 0.988 & 0.892 \\
时延与噪声耦合 & 0.93 & 0.039 & 0.971 & 0.842 \\
\bottomrule
\end{tabular}
\end{table}

从表\ref{tab:ch6_robustness}可见，即使在耦合扰动场景下，重规划成功率仍保持在 $97\%$ 以上，说明补偿机制能够有效控制误差传播。需要指出的是，强扰动条件下轨迹误差与碰撞概率均有上升，提示系统仍需在极端通信退化场景中引入更保守的安全边界。

\begin{algorithm}[H]
\caption{在线增量重规划与容错执行伪代码}
\label{alg:ch6_replan}
\begin{algorithmic}[1]
\REQUIRE 当前状态 $\mathbf{x}_k$，全局路径 $\mathcal{P}_g$，扰动估计 $(\Delta t_k,\mathbf{n}_k)$
\ENSURE 控制动作 $\mathbf{u}_k$
\STATE 计算补偿目标 $\hat{\mathbf{g}}_k$（式\eqref{eq:ch6_delay_compensation_model}）
\STATE 根据式\eqref{eq:ch6_switch_rule}评估策略切换条件
\IF{$\mathcal{S}_k>\tau_{sw}$}
    \STATE 由 D3QN 输出局部动作候选
    \STATE 执行动态安全边界检查
\ELSE
    \STATE 采用 A* 引导动作并执行平滑校正
\ENDIF
\IF{轨迹误差超过阈值或局部碰撞风险升高}
    \STATE 触发增量重规划并更新局部参考段
\ENDIF
\STATE 输出 $\mathbf{u}_k$ 并记录诊断信息
\end{algorithmic}
\end{algorithm}

算法\ref{alg:ch6_replan}强调“补偿-切换-校核-重规划”的时序闭环。该机制将故障检测与决策执行统一到同一周期内，减少了传统串行架构中“先失败再修复”的响应迟滞。

\subsection{多任务长航程综合评估与工程可行性}
\label{subsec:ch6_long_horizon}

为评估算法在长期任务中的稳定收益，本节构建多任务长航程实验，包含“单目标快速到达”“多目标巡检”“动态风险穿越”三类任务。定义综合任务效用为
\begin{equation}
\mathcal{U}_{mission}=w_l\frac{L_{ref}}{L}+w_s(1-P_{col})+w_t\frac{T_{ref}}{T}+w_e\frac{E_{ref}}{E},
\label{eq:ch6_mission_utility}
\end{equation}
其中 $L$、$T$、$E$ 分别表示路径长度、任务时长与能耗，$P_{col}$ 表示碰撞概率。

\begin{table}[H]
\centering
\small
\caption{多任务长航程场景下的综合性能对比}
\label{tab:ch6_long_horizon}
\begin{tabular}{p{0.20\textwidth}p{0.14\textwidth}p{0.14\textwidth}p{0.14\textwidth}p{0.14\textwidth}}
\toprule
任务类型 & 路径长度(m) & 任务时长(s) & 碰撞概率 & $\mathcal{U}_{mission}$ \\
\midrule
单目标快速到达 & 652.7 & 684.2 & 0.010 & 0.921 \\
多目标巡检任务 & 1014.5 & 1098.3 & 0.014 & 0.903 \\
动态风险穿越 & 735.8 & 812.7 & 0.021 & 0.884 \\
多目标+动态风险耦合 & 1126.9 & 1215.4 & 0.028 & 0.861 \\
\bottomrule
\end{tabular}
\end{table}

结果表明，随着任务复杂度提升，综合效用呈可解释下降趋势，但整体仍保持在较高水平。特别是在“多目标+动态风险耦合”场景中，算法未出现不可恢复失效，说明所提混合规划框架具备面向长航程任务的工程可行性。

\begin{figure}[H]
\centering
\resizebox{0.70\textwidth}{!}{%
\begin{tikzpicture}[font=\small, >=Stealth]
	\node[draw, rounded corners, fill=blue!10, align=center, minimum width=7.1cm, minimum height=1.0cm] (m1) at (5,8.6) {任务队列生成与优先级排序};
	\node[draw, rounded corners, fill=cyan!12, align=center, minimum width=7.1cm, minimum height=1.0cm] (m2) at (5,6.8) {全局可达性分析与子目标分配};
	\node[draw, rounded corners, fill=green!12, align=center, minimum width=7.1cm, minimum height=1.0cm] (m3) at (5,5.0) {A*-D3QN 混合规划执行};
	\node[draw, rounded corners, fill=orange!14, align=center, minimum width=7.1cm, minimum height=1.0cm] (m4) at (5,3.2) {安全核验与在线重规划触发};
	\node[draw, rounded corners, fill=purple!10, align=center, minimum width=7.1cm, minimum height=1.0cm] (m5) at (5,1.4) {任务完成评估与策略参数回写};
	\draw[->, thick, draw=black!75] (m1) -- (m2);
	\draw[->, thick, draw=black!75] (m2) -- (m3);
	\draw[->, thick, draw=black!75] (m3) -- (m4);
	\draw[->, thick, draw=black!75] (m4) -- (m5);
	\draw[->, thick, draw=blue!70!black] (m5.west) .. controls (1.2,1.2) and (1.2,7.8) .. node[left, align=center]{阶段复盘\\与调参} (m2.west);
\end{tikzpicture}%
}
\caption{多任务长航程规划执行程序框图}
\label{fig:ch6_long_horizon_flow}
\end{figure}

图\ref{fig:ch6_long_horizon_flow}体现了长航程任务中“任务管理-规划执行-安全核验-策略回写”的闭环组织方式。相比单任务场景，该流程增加了任务级调度与阶段复盘环节，使规划器能够在长时序运行中保持稳定性能。

% === 6.5_本章小结.tex ===
\section{本章小结}
\label{sec:ch6_summary}

本章围绕“数字孪生约束下的月面巡视器路径规划”展开研究，完成了从问题建模、混合算法设计到仿真/地面验证的完整闭环。
主要贡献如下：
\begin{myitemize}
\item 构建了环境建模、状态感知、决策优化、虚实校核一体化规划框架，实现了规划链路与数字孪生系统接口的工程对齐；
\item 提出 A* 搜索算法（A-star）与深度双对偶Q网络（D3QN）协同的混合规划方法，并通过改进 A* 平滑策略、D3QN 轻量推理与策略切换机制实现全局可达性与局部响应性的兼顾；
\item 通过多复杂度仿真和地面平台验证，定量证明了方法在路径平稳性、碰撞抑制与综合效率方面的优势。
\end{myitemize}

与现有单一全局或单一强化学习方法相比，本章方法的技术优势在于：
将“全局初始约束”与“局部学习决策”统一到同一闭环框架中，避免了纯学习方法可达性不足和纯搜索方法动态适应性不足的问题。

研究局限在于：
当前地面可视化端帧率仍低于工程目标，且部分时延补偿结论基于控制仿真模型而非全实物在环平台。
后续工作将面向三方面展开：
\begin{myitemize}
\item 引入增量式全局重规划与局部策略并行执行，降低复杂场景下秒级规划开销；
\item 开展深度双对偶Q网络（D3QN）模型压缩与推理加速，提升端侧部署效率；
\item 构建更高保真地面在环试验链路，系统评估通讯时延、感知误差与执行偏差的耦合影响。
\end{myitemize}


\end{document}
